% Equations et inéquations su 1er degré à une inconnue dans IR — Mathématiques 3ème — Cours signé OnBuch+
% Programme officiel MINESEC. Compiler avec : tectonic N09-equations-inequations-su-1er-degre-inconnue-ir.tex
\newcommand{\DOCMATIERE}{Mathématiques}
\newcommand{\DOCNIVEAU}{3ème}
\newcommand{\DOCMODULE}{Mathématiques – classe de 3ème}
\newcommand{\DOCLECON}{N09}
\newcommand{\DOCTITRE}{Equations et inéquations su 1er degré à une inconnue dans IR}
% OnBuch+ — préambule « cours signés » ENRICHI (compilé avec tectonic / XeLaTeX)
% Système visuel « Pop » d'OnBuch+ : fond crème (jamais blanc), bordures encre
% épaisses, ombres « dures » décalées, titres Archivo Black en capitales.
% Version MATHÉMATIQUES : graphes (pgfplots/tikz), boîtes pédagogiques
% (définition, propriété, méthode, attention, le savais-tu, exercice résolu,
% exercice, TP, bilan), page de garde et signature OnBuch+ ancrée en bas.
%
% Macros attendues dans main.tex AVANT \input{preamble} :
%   \DOCMATIERE  \DOCNIVEAU  \DOCMODULE  \DOCLECON (numéro)  \DOCTITRE
\documentclass[11pt]{article}

\usepackage{fontspec}
\usepackage[french]{babel}
\usepackage[a4paper,margin=2cm,top=3.4cm,bottom=2.8cm,headheight=1.4cm,footskip=1.3cm]{geometry}
\usepackage{amsmath,amssymb,amsfonts}
\usepackage{mathtools} % \xmapsto, \coloneqq... souvent utilisés par les modèles
\usepackage{cancel} % simplifications barrées
% \abs{z} (module, valeur absolue) et \norm{u} (norme), utilisés par les modèles
\providecommand{\abs}{}\let\abs\relax\DeclarePairedDelimiter{\abs}{\lvert}{\rvert}
\providecommand{\norm}{}\let\norm\relax\DeclarePairedDelimiter{\norm}{\lVert}{\rVert}
\usepackage[table]{xcolor} % [table] : fournit aussi \rowcolors (alternance de lignes)
\usepackage{pagecolor}
\usepackage{tikz}
\usetikzlibrary{arrows.meta,positioning,calc,shapes.geometric,shapes.misc,shapes.arrows,decorations.pathmorphing,decorations.pathreplacing,decorations.markings,patterns,shadows,mindmap,trees,fit,backgrounds,matrix,intersections,angles,quotes}
\usepackage{pgfplots}
\usepackage{pgf-pie} % diagrammes circulaires \pie (statistiques)
\pgfplotsset{compat=1.18}
\usepgfplotslibrary{fillbetween,groupplots} % aires entre courbes (calcul intégral)
\usepackage{enumitem}
\usepackage{fancyhdr}
% [most] charge toutes les bibliothèques tcolorbox (listings, minted, raster,
% documentation...) et gonfle inutilement la mémoire TeX sur un document long
% (~300 boîtes) : on ne charge que ce qui est réellement utilisé.
\usepackage[skins,breakable]{tcolorbox}
\usepackage{graphicx}
\usepackage{colortbl}
\usepackage{booktabs}
\usepackage{tabularx}
\usepackage{diagbox} % case d'en-tête diagonale des tableaux à double entrée
% Colonnes extensibles centrée / gauche / droite (C, L, R), souvent utilisées par les modèles
\newcolumntype{C}{>{\centering\arraybackslash}X}
\newcolumntype{L}{>{\raggedright\arraybackslash}X}
\newcolumntype{R}{>{\raggedleft\arraybackslash}X}
\usepackage{array}
\usepackage{multicol}
\usepackage{siunitx}
% parse-numbers=false : accepte aussi \SI{2,5 \times 10^{-3}}{...}
\sisetup{locale=FR,output-decimal-marker={,},group-minimum-digits=4,parse-numbers=false}
\DeclareSIUnit\ohm{\ensuremath{\symup{\Omega}}} % U+2126 absent de la police
\usepackage{qrcode}
% \degree : commande fréquemment utilisée par les modèles pour un angle ou une
% température en dehors de \SI{}{} (non fournie par les paquets chargés).
\providecommand{\degree}{^{\circ}}
% \qed (fin de démonstration), utilisé par les modèles sans amsthm
\providecommand{\qed}{\hfill$\square$}
% \barre : double barre des valeurs interdites dans un tableau de variations (tkz-tab)
\providecommand{\barre}{\ensuremath{\|}}
% environnement proof (amsthm non chargé) utilisé par les modèles
\ifdefined\proof\else\newenvironment{proof}[1][Démonstration]{\par\noindent\textit{#1.}\ }{\qed\par}\fi
\usepackage{lastpage}
\usepackage{hyperref}
\hypersetup{hidelinks}

% ============================================================
% Palette « Pop » OnBuch+ — mêmes hex que la classe P (plus_ui.dart)
% ============================================================
\definecolor{poporange}{HTML}{F59321}
\definecolor{popdark}{HTML}{E07A0C}
\definecolor{popcream}{HTML}{FFF6E8}
\definecolor{popink}{HTML}{1C1714}
\definecolor{poppurple}{HTML}{7A5AE0}
\definecolor{popgreen}{HTML}{2E9E6A}
\definecolor{poppink}{HTML}{E0457B}
\definecolor{popblue}{HTML}{2F7DE1}
\definecolor{popgold}{HTML}{C98A12}
\definecolor{popmuted}{HTML}{8A7F76}
\definecolor{popline}{HTML}{EADFD2}
\definecolor{popgray}{HTML}{8A7F76}
\colorlet{popgrey}{popgray}
% Alias tolérants (le modèle nomme parfois les couleurs différemment)
\colorlet{popwhite}{white}
\colorlet{popblack}{popink}
\colorlet{popred}{poppink}
\colorlet{popyellow}{popgold}
% Teintes claires pour fonds de tableaux / zones
\colorlet{popblueL}{popblue!10!white}
\colorlet{poporangeL}{poporange!14!white}
\colorlet{popgreenL}{popgreen!12!white}
\colorlet{poppurpleL}{poppurple!12!white}
\colorlet{poppinkL}{poppink!12!white}
\colorlet{popgoldL}{popgold!14!white}

\pagecolor{popcream}

% ============================================================
% Typographie
% ============================================================
\defaultfontfeatures{Ligatures=TeX}
\setmainfont{PlusJakartaSans-Medium.ttf}[Path=../../../fonts/,BoldFont=PlusJakartaSans-Bold.ttf]
% Maths en sans-serif (Fira Math) avec chiffres et lettres droites de Plus
% Jakarta, pour que les formules se fondent dans le texte.
\usepackage[math-style=ISO,bold-style=ISO]{unicode-math}
\setmathfont{FiraMath-Regular.otf}[Path=../../../fonts/]
\setmathfont{PlusJakartaSans-Medium.ttf}[Path=../../../fonts/,range={up/num,up/{Latin,latin},\mathup}]
\usepackage{icomma} % virgule décimale sans espace en mode math
% Caractères mathématiques Unicode absents des polices de texte (Plus Jakarta,
% Archivo Black) : rendus par leur équivalent mathématique, en texte comme en
% formule — sinon ils s'affichent en carré vide (ex. « Arithmétique dans ℤ »).
\usepackage{newunicodechar}
\newunicodechar{ℝ}{\ensuremath{\mathbb{R}}}
\newunicodechar{ℤ}{\ensuremath{\mathbb{Z}}}
\newunicodechar{ℂ}{\ensuremath{\mathbb{C}}}
\newunicodechar{ℕ}{\ensuremath{\mathbb{N}}}
\newunicodechar{ℚ}{\ensuremath{\mathbb{Q}}}
\newunicodechar{𝔻}{\ensuremath{\mathbb{D}}}
\newunicodechar{α}{\ensuremath{\alpha}}
\newunicodechar{β}{\ensuremath{\beta}}
\newunicodechar{γ}{\ensuremath{\gamma}}
\newunicodechar{δ}{\ensuremath{\delta}}
\newunicodechar{ε}{\ensuremath{\varepsilon}}
\newunicodechar{θ}{\ensuremath{\theta}}
\newunicodechar{λ}{\ensuremath{\lambda}}
\newunicodechar{μ}{\ensuremath{\mu}}
\newunicodechar{σ}{\ensuremath{\sigma}}
\newunicodechar{τ}{\ensuremath{\tau}}
\newunicodechar{φ}{\ensuremath{\varphi}}
\newunicodechar{ψ}{\ensuremath{\psi}}
\newunicodechar{ω}{\ensuremath{\omega}}
\newunicodechar{Σ}{\ensuremath{\Sigma}}
% majuscules produites par \MakeUppercase dans les titres (α-aminés -> Α)
\newunicodechar{Α}{\ensuremath{\alpha}}
\newunicodechar{Β}{\ensuremath{\beta}}
\newunicodechar{Γ}{\ensuremath{\Gamma}}
\newunicodechar{Δ}{\ensuremath{\Delta}}
\newunicodechar{Ε}{\ensuremath{\varepsilon}}
\newunicodechar{Θ}{\ensuremath{\Theta}}
\newunicodechar{Λ}{\ensuremath{\Lambda}}
\newunicodechar{Μ}{\ensuremath{\mu}}
\newunicodechar{Τ}{\ensuremath{\tau}}
\newunicodechar{Φ}{\ensuremath{\Phi}}
\newunicodechar{Ψ}{\ensuremath{\Psi}}
\newunicodechar{Ω}{\ensuremath{\Omega}}
\newunicodechar{↦}{\ensuremath{\mapsto}}
\newunicodechar{∈}{\ensuremath{\in}}
\newunicodechar{∉}{\ensuremath{\notin}}
\newunicodechar{∧}{\ensuremath{\wedge}}
\newunicodechar{⇔}{\ensuremath{\Leftrightarrow}}
\newunicodechar{⟺}{\ensuremath{\Longleftrightarrow}}
\newunicodechar{⇒}{\ensuremath{\Rightarrow}}
\newunicodechar{⟹}{\ensuremath{\Longrightarrow}}
\newunicodechar{∘}{\ensuremath{\circ}}
\newunicodechar{∪}{\ensuremath{\cup}}
\newunicodechar{∩}{\ensuremath{\cap}}
\newunicodechar{≡}{\ensuremath{\equiv}}
\newunicodechar{⊂}{\ensuremath{\subset}}
\newunicodechar{⊥}{\ensuremath{\perp}}
\newunicodechar{∃}{\ensuremath{\exists}}
\newunicodechar{∀}{\ensuremath{\forall}}
\newunicodechar{∛}{\ensuremath{\sqrt[3]{\,}}}
\newunicodechar{∜}{\ensuremath{\sqrt[4]{\,}}}
\newunicodechar{⃗}{\ensuremath{{}^{\to}}}
\newunicodechar{ⁿ}{\ifmmode^{n}\else\textsuperscript{\ensuremath{n}}\fi}
\newunicodechar{ˣ}{\ifmmode^{x}\else\textsuperscript{\ensuremath{x}}\fi}
\newunicodechar{ᵇ}{\ifmmode^{b}\else\textsuperscript{\ensuremath{b}}\fi}
\newunicodechar{ʳ}{\ifmmode^{r}\else\textsuperscript{\ensuremath{r}}\fi}
\newunicodechar{ᵘ}{\ifmmode^{u}\else\textsuperscript{\ensuremath{u}}\fi}
\newunicodechar{ʸ}{\ifmmode^{y}\else\textsuperscript{\ensuremath{y}}\fi}
\newunicodechar{ᵃ}{\ifmmode^{a}\else\textsuperscript{\ensuremath{a}}\fi}
\newunicodechar{ᵉ}{\ifmmode^{e}\else\textsuperscript{\ensuremath{e}}\fi}
\newunicodechar{ᵏ}{\ifmmode^{k}\else\textsuperscript{\ensuremath{k}}\fi}
\newunicodechar{ᵖ}{\ifmmode^{p}\else\textsuperscript{\ensuremath{p}}\fi}
\newunicodechar{ᴺ}{\ifmmode^{N}\else\textsuperscript{\ensuremath{N}}\fi}
\newunicodechar{ᶜ}{\ifmmode^{c}\else\textsuperscript{\ensuremath{c}}\fi}
\newunicodechar{ʰ}{\ifmmode^{h}\else\textsuperscript{\ensuremath{h}}\fi}
\newunicodechar{ᵠ}{\ifmmode^{q}\else\textsuperscript{\ensuremath{q}}\fi}
\newunicodechar{ₙ}{\ifmmode_{n}\else\textsubscript{\ensuremath{n}}\fi}
\newunicodechar{ₐ}{\ifmmode_{a}\else\textsubscript{\ensuremath{a}}\fi}
\newunicodechar{ᵢ}{\ifmmode_{i}\else\textsubscript{\ensuremath{i}}\fi}
\newunicodechar{ᵣ}{\ifmmode_{r}\else\textsubscript{\ensuremath{r}}\fi}
\newunicodechar{ₖ}{\ifmmode_{k}\else\textsubscript{\ensuremath{k}}\fi}
\newunicodechar{ᵦ}{\ifmmode_{\beta}\else\textsubscript{\ensuremath{\beta}}\fi}
\newunicodechar{ₓ}{\ifmmode_{x}\else\textsubscript{\ensuremath{x}}\fi}
\newfontfamily\popdisplay{ArchivoBlack-Regular.ttf}[Path=../../../fonts/]
\newfontfamily\poplogo{SpaceGrotesk-Bold.ttf}[Path=../../../fonts/]
\newfontfamily\popsemi{PlusJakartaSans-SemiBold.ttf}[Path=../../../fonts/]
\newfontfamily\popxbold{PlusJakartaSans-ExtraBold.ttf}[Path=../../../fonts/]
\newfontfamily\popxboldls{PlusJakartaSans-ExtraBold.ttf}[Path=../../../fonts/,LetterSpace=3.0]

\setlist[enumerate]{leftmargin=1.7em,itemsep=5pt,topsep=4pt}
\setlist[itemize]{leftmargin=1.7em,itemsep=4pt,topsep=4pt,label={\color{poporange}\textbullet}}
\setlength{\parindent}{0pt}
\setlength{\parskip}{5pt}
\renewcommand{\arraystretch}{1.25}

% pgfplots : style Pop par défaut
\pgfplotsset{
  every axis/.append style={axis line style={popink,line width=1pt},tick style={popink},
    grid style={popline,line width=0.5pt},label style={font=\small},tick label style={font=\footnotesize}},
  popcurve/.style={poporange,line width=1.8pt,no markers,smooth},
}
% Même style disponible dans un \draw TikZ simple (hors pgfplots)
\tikzset{popcurve/.style={poporange,line width=1.8pt}}

% ============================================================
% Pill / squiggle
% ============================================================
\newcommand{\poppill}[3]{%
  \tikz[baseline=-0.55ex]\node[draw=popink,line width=1.2pt,rounded corners=2.6pt,%
    fill=#2,inner xsep=6.5pt,inner ysep=3pt]{%
    \popxboldls\fontsize{7}{7}\selectfont\color{#3}\MakeUppercase{#1}};%
}
\newcommand{\popsquiggle}[1]{%
  \tikz[baseline=-0.4ex]{
    \draw[#1,line width=2.6pt,line cap=round]
      plot [smooth] coordinates {(0,0.09) (0.34,0.01) (0.68,0.21) (1.02,0.01) (1.36,0.10)};
  }%
}
% Mot-clé surligné (marqueur orange sous le texte)
\newcommand{\cle}[1]{{\popxbold\color{popdark}#1}}

% ============================================================
% En-tête / pied
% ============================================================
\pagestyle{fancy}
\fancyhf{}
\renewcommand{\headrulewidth}{0pt}
\renewcommand{\footrulewidth}{0pt}
\lhead{\raisebox{-0.15ex}{\poplogo\fontsize{13}{13}\selectfont\color{popink}OnBuch\color{poporange}+}}
\rhead{\poppill{\DOCMATIERE\ \textbullet\ \DOCNIVEAU\ \textbullet\ Leçon \DOCLECON}{white}{popink}}
\lfoot{\popxbold\fontsize{7.6}{7.6}\selectfont\color{popmuted}\MakeUppercase{Cours signé OnBuch+}\ \textbullet\ {\popsemi\DOCTITRE}}
\rfoot{\tikz[baseline=-0.6ex]\node[rounded corners=8.5pt,fill=popink,minimum height=17pt,minimum width=17pt,inner xsep=5pt,inner sep=0]{\popxbold\fontsize{8}{8}\selectfont\color{popcream}\thepage\,/\,\pageref*{LastPage}};}

% ============================================================
% Titres
% ============================================================
\newcommand{\chaptitre}[2]{%
  \begin{tcolorbox}[enhanced,colback=poporange,colframe=popink,boxrule=2.6pt,arc=11pt,
      drop shadow=popink,left=15pt,right=15pt,top=12pt,bottom=11pt,before skip=3pt,after skip=18pt]
    \poppill{#2}{popink}{popcream}\par\vspace{9pt}
    {\raggedright\hyphenpenalty=10000\exhyphenpenalty=10000\popdisplay\fontsize{22}{24}\selectfont\color{popcream}\MakeUppercase{#1}\par}
  \end{tcolorbox}
}
% Section numérotée : pastille orange avec le numéro + titre Archivo
\newcounter{popsec}
\newcommand{\coursec}[1]{%
  \stepcounter{popsec}\setcounter{popsub}{0}%
  \par\vspace{16pt}\noindent
  \begin{minipage}{\linewidth}
  \tikz[baseline=-0.7ex]\node[draw=popink,line width=1.6pt,fill=poporange,rounded corners=4pt,minimum size=22pt,inner sep=0]{\popdisplay\fontsize{12}{12}\selectfont\color{popcream}\thepopsec};%
  \hspace{8pt}{\popdisplay\fontsize{15}{17}\selectfont\color{popink}\MakeUppercase{#1}}\par
  \vspace{4pt}\noindent{\color{poporange}\rule{36pt}{3.6pt}}
  \end{minipage}\par\nopagebreak\vspace{8pt}
}
\newcounter{popsub}
\newcommand{\courssub}[1]{%
  \stepcounter{popsub}%
  \par\vspace{9pt}\noindent{\popxbold\fontsize{11.5}{13}\selectfont\color{popink}{\color{poporange}\thepopsec.\thepopsub}\hspace{6pt}#1}\par\nopagebreak\vspace{4pt}
}

% ============================================================
% Boîtes pédagogiques — même recette : fond blanc, bordure épaisse colorée,
% ombre dure encre, badge plein. Titre optionnel [#1].
% ============================================================
\tcbset{popbase/.style={breakable,enhanced,colback=white,boxrule=2.3pt,arc=9pt,
  shadow={3pt}{-3pt}{0pt}{popink},left=14pt,right=14pt,top=11pt,bottom=11pt,before skip=11pt,after skip=15pt,
  fontupper=\popsemi\fontsize{10.6}{14.5}\selectfont\color{popink},
  fontlower=\popsemi\fontsize{10.6}{14.5}\selectfont\color{popink}}}
% Coche dessinée (le glyphe ✓ n'existe pas dans les polices du document)
\newcommand{\popcheck}[1]{\tikz[baseline=-0.5ex]\draw[#1,line width=1.6pt,line cap=round,line join=round](-0.13,0.0)--(-0.04,-0.09)--(0.13,0.1);}
\newcommand{\popboxhead}[3]{\poppill{#1}{#2}{white}\ifx\relax#3\relax\else\hspace{6pt}{\popxbold\fontsize{10.6}{12}\selectfont\color{popink}#3}\fi\par\vspace{7pt}}

\newtcolorbox{aretenir}[1][]{popbase,colframe=popgreen,before upper={\popboxhead{À retenir}{popgreen}{#1}}}
\newtcolorbox{exemplebox}[1][]{popbase,colframe=poporange,before upper={\popboxhead{Exemple}{poporange}{#1}}}
\newtcolorbox{definition}[1][]{popbase,colframe=popblue,before upper={\popboxhead{Définition}{popblue}{#1}}}
\newtcolorbox{propriete}[1][]{popbase,colframe=poppurple,before upper={\popboxhead{Propriété}{poppurple}{#1}}}
\newtcolorbox{methode}[1][]{popbase,colframe=popgold,before upper={\popboxhead{Méthode}{popgold}{#1}}}
\newtcolorbox{attention}[1][]{popbase,colframe=poppink,before upper={\popboxhead{Attention}{poppink}{#1}}}
\newtcolorbox{experience}[1][]{popbase,colframe=popblue,colback=popblueL,before upper={\popboxhead{Expérience}{popblue}{#1}}}
% « Le savais-tu ? » : carte violette pleine (contexte, culture, Cameroun)
\newtcolorbox{savaistu}[1][]{popbase,colback=poppurple,colframe=popink,
  fontupper=\popsemi\fontsize{10.4}{14.2}\selectfont\color{white},
  before upper={\poppill{Le savais-tu\,?}{white}{poppurple}\ifx\relax#1\relax\else\hspace{6pt}{\popxbold\color{white}#1}\fi\par\vspace{7pt}}}
% Exercice résolu : énoncé en haut, \tcblower puis la solution sur fond crème
\newcounter{popexo}\newcounter{popexr}
\newtcolorbox{exoresolu}[1][]{popbase,colframe=poporange,colbacklower=popcream,
  segmentation style={popink,line width=1.2pt,dashed},
  before upper={\stepcounter{popexr}\popboxhead{Exercice résolu \thepopexr}{poporange}{#1}},
  before lower={\poppill{Solution}{popink}{popcream}\par\vspace{6pt}}}
% Exercice (non résolu en place) — la correction est dans la partie Corrigés
\newtcolorbox{exercice}[1][]{popbase,colframe=popink,
  before upper={\stepcounter{popexo}\popboxhead{Exercice \thepopexo}{popink}{#1}}}
% Corrigé
\newtcolorbox{corrige}[1][]{popbase,colframe=popgreen,colback=popgreenL,
  before upper={\popboxhead{Corrigé}{popgreen}{#1}}}
% Objectifs / prérequis / bilan
\newtcolorbox{objectifs}{popbase,colframe=popink,colback=poporangeL,
  before upper={\popboxhead{Objectifs de la leçon}{popink}{}}}
\newtcolorbox{prerequis}{popbase,colframe=popink,colback=popblueL,
  before upper={\popboxhead{Prérequis}{popblue}{}}}
\newtcolorbox{bilan}{popbase,colframe=popink,colback=white,boxrule=2.8pt,
  before upper={\popboxhead{Fiche bilan}{poporange}{L'essentiel à connaître pour l'examen}}}
% Figure encadrée avec légende
\newtcolorbox{popfigure}[1][]{enhanced,breakable=false,colback=white,colframe=popline,boxrule=1.4pt,arc=8pt,
  left=10pt,right=10pt,top=10pt,bottom=8pt,before skip=10pt,after skip=12pt,halign=center,
  lower separated=false,halign lower=center,
  fontlower=\popsemi\fontsize{9}{11.5}\selectfont\color{popmuted},
  #1}
\newcommand{\legende}[1]{\par\vspace{6pt}{\popsemi\fontsize{9}{11.5}\selectfont\color{popmuted}#1}}

% ============================================================
% Signature OnBuch+
% ============================================================
\newcommand{\poplogoinline}[1]{{\poplogo\fontsize{#1}{#1}\selectfont\color{popink}OnBuch\color{poporange}+}}
\newcommand{\signatureonbuch}{%
  \begin{tcolorbox}[enhanced,colback=popink,colframe=popink,boxrule=2.2pt,arc=10pt,
      drop shadow=poporange,left=14pt,right=14pt,top=10pt,bottom=10pt,before skip=0pt,after skip=0pt]
    \tikz[baseline=-0.7ex]\node[circle,fill=popgreen,draw=popcream,line width=1.3pt,minimum size=22pt,inner sep=0]{\popcheck{white}};%
    \hspace{9pt}\begin{minipage}[c]{0.62\linewidth}
      {\popxbold\fontsize{12}{13}\selectfont\color{popcream}Signé\ \textbullet\ {\poplogo OnBuch\color{poporange}+}}\par\vspace{2pt}
      {\raggedright\popsemi\fontsize{8}{10.5}\selectfont\color{popline}\MakeUppercase{Équipe pédagogique OnBuch+}\\\MakeUppercase{Conforme au programme officiel MINESEC}\par}
    \end{minipage}\hfill
    {\poplogo\fontsize{22}{22}\selectfont\color{popcream}OnBuch\color{poporange}+}
  \end{tcolorbox}%
}

% ============================================================
% Page de garde
% #1 titre · #2 matière · #3 niveau · #4 module · #5 n° leçon · #6 durée ·
% #7 description / objectifs (texte ou itemize)
% ============================================================
\newcommand{\pagegarde}[7]{%
  \thispagestyle{empty}
  \newgeometry{a4paper,margin=1.7cm,top=1.6cm,bottom=1.4cm}%
  \begin{tikzpicture}[remember picture,overlay]
    \fill[poporange!18] ([xshift=-3.2cm,yshift=-3.6cm]current page.north east) circle (4.6cm);
    \fill[poppurple!14] ([xshift=2.4cm,yshift=7.6cm]current page.south west) circle (3.8cm);
  \end{tikzpicture}%
  {\poplogo\fontsize{30}{30}\selectfont\color{popink}OnBuch\color{poporange}+}\hfill
  \raisebox{0.6ex}{\poppill{Cours signé \textbullet\ Premium}{popink}{popcream}}\par
  \vspace{1pt}\popsquiggle{poppurple}\par
  \vspace{8pt}
  \begin{tcolorbox}[enhanced,colback=poporange,colframe=popink,boxrule=2.8pt,arc=13pt,
      drop shadow=popink,left=18pt,right=18pt,top=12pt,bottom=13pt,after skip=10pt]
    \poppill{#2 \textbullet\ #3}{popink}{popcream}\hspace{5pt}\poppill{#4}{white}{popink}\par\vspace{8pt}
    {\popdisplay\fontsize{14}{14}\selectfont\color{popink}LEÇON #5}\par\vspace{4pt}
    {\raggedright\hyphenpenalty=10000\exhyphenpenalty=10000\popdisplay\fontsize{25}{27}\selectfont\color{popcream}\MakeUppercase{#1}\par}
    \vspace{8pt}
    {\popxbold\fontsize{9.5}{9.5}\selectfont\color{popink}\MakeUppercase{Durée conseillée : #6}}
  \end{tcolorbox}
  \begin{tcolorbox}[enhanced,colback=white,colframe=popink,boxrule=2.2pt,arc=10pt,
      drop shadow=popink,left=16pt,right=16pt,top=10pt,bottom=10pt,after skip=8pt]
    \poppill{Dans cette leçon}{poppurple}{white}\par\vspace{5pt}
    \popsemi\fontsize{9.3}{12.6}\selectfont\color{popink}#7
  \end{tcolorbox}
  \noindent\begin{minipage}[t]{0.487\linewidth}
    \begin{tcolorbox}[enhanced,colback=popgreen,colframe=popink,boxrule=2.2pt,arc=10pt,
        drop shadow=popink,left=11pt,right=11pt,top=10pt,bottom=10pt,height=3.1cm,valign=center]
      \tcbox[colback=white,colframe=popink,boxrule=1.3pt,arc=5pt,boxsep=0pt,nobeforeafter,
        left=3pt,right=3pt,top=3pt,bottom=3pt,valign=center]{\qrcode[height=1.9cm]{https://play.google.com/store/apps/details?id=com.luvvix.onbuch&hl=fr}}%
      \hspace{8pt}\begin{minipage}[c]{0.5\linewidth}
        {\popdisplay\fontsize{11}{12.5}\selectfont\color{white}\MakeUppercase{Télécharge OnBuch}}\par\vspace{4pt}
        {\raggedright\popsemi\fontsize{8.4}{11}\selectfont\color{white}Gratuit \textbullet\ Google Play\\Tuteur IA, annales, concours, résultats\par}
      \end{minipage}
    \end{tcolorbox}
  \end{minipage}\hfill
  \begin{minipage}[t]{0.487\linewidth}
    \begin{tcolorbox}[enhanced,colback=poppurple,colframe=popink,boxrule=2.2pt,arc=10pt,
        drop shadow=popink,left=11pt,right=11pt,top=10pt,bottom=10pt,height=3.1cm,valign=center]
      \tikz[baseline=-0.7ex]\node[circle,fill=white,draw=popink,line width=1.3pt,minimum size=1.6cm,inner sep=0]{\popdisplay\fontsize{24}{24}\selectfont\color{poppurple}+};%
      \hspace{8pt}\begin{minipage}[c]{0.6\linewidth}
        {\popdisplay\fontsize{11}{12.5}\selectfont\color{white}\MakeUppercase{Passe à OnBuch+}}\par\vspace{4pt}
        {\raggedright\popsemi\fontsize{8.4}{11}\selectfont\color{white}Tous les cours signés, Tuteur IA illimité, fiches PDF — onglet OnBuch+ de l'app\par}
      \end{minipage}
    \end{tcolorbox}
  \end{minipage}
  \vspace{8pt}
  \popcontenu
  \vfill
  \signatureonbuch
  \restoregeometry
  \newpage
}

% Rangée « Ce cours contient » (page de garde)
\newcommand{\poptile}[3]{% #1 couleur #2 gros texte #3 légende
  \begin{tcolorbox}[enhanced,colback=#1,colframe=popink,boxrule=2pt,arc=9pt,shadow={2.5pt}{-2.5pt}{0pt}{popink},
      left=6pt,right=6pt,top=9pt,bottom=9pt,halign=center,valign=center,height=1.75cm,width=\linewidth,nobeforeafter]
    {\popdisplay\fontsize{12.5}{13}\selectfont\color{white}\MakeUppercase{#2}}\par\vspace{3pt}
    {\centering\popsemi\fontsize{7.8}{9.5}\selectfont\color{white}#3\par}
  \end{tcolorbox}}
\newcommand{\popcontenu}{%
  \par\noindent{\popxboldls\fontsize{7.5}{7.5}\selectfont\color{popmuted}CE COURS SIGNÉ CONTIENT}\par\vspace{7pt}
  \noindent\begin{minipage}[t]{0.235\linewidth}\poptile{popblue}{Cours}{complet, illustré, pas à pas}\end{minipage}\hfill
  \begin{minipage}[t]{0.235\linewidth}\poptile{poporange}{Méthodes}{et exercices résolus}\end{minipage}\hfill
  \begin{minipage}[t]{0.235\linewidth}\poptile{poppink}{Exercices}{type examen + corrigés}\end{minipage}\hfill
  \begin{minipage}[t]{0.235\linewidth}\poptile{popgreen}{Fiche bilan}{l'essentiel à retenir}\end{minipage}\par}

% Sommaire
\newtcolorbox{sommaire}{enhanced,colback=white,colframe=popink,boxrule=2.2pt,arc=10pt,
  drop shadow=popink,left=18pt,right=18pt,top=13pt,bottom=13pt,before skip=6pt,after skip=20pt,
  before upper={\poppill{Sommaire}{poppurple}{white}\par\vspace{9pt}\popsemi\fontsize{10.6}{14.5}\selectfont\color{popink}}}

% Signature de fin de cours — poussée tout en bas de la dernière page
\newcommand{\findecours}{%
  \par\vfill\nopagebreak
  \begin{center}\popsquiggle{poporange}\end{center}\vspace{4pt}
  \signatureonbuch
}
\AtBeginDocument{\def\llbracket{\mathopen{[\mkern-3.5mu[}}\def\rrbracket{\mathclose{]\mkern-3.5mu]}}} % intervalles d'entiers (glyphe ⟦ absent de la police)
% Glyphes absents de Fira Math : redéfinis à partir de glyphes disponibles
\AtBeginDocument{%
  \def\setminus{\mathbin{\backslash}}\let\smallsetminus\setminus
  \def\vdots{\mathord{\vbox{\baselineskip3.2pt\lineskiplimit0pt\kern4pt\hbox{.}\hbox{.}\hbox{.}}}}%
  \def\ddots{\mathinner{\mkern1mu\raise6.5pt\hbox{.}\mkern2mu\raise3.6pt\hbox{.}\mkern2mu\raise0.7pt\hbox{.}\mkern1mu}}%
  \def\triangle{\mathord{\scalebox{0.9}{$\symup{\Delta}$}}}%
  \def\nabla{\mathord{\raisebox{\depth}{\scalebox{1}[-1]{$\symup{\Delta}$}}}}%
  \def\checkmark{\popcheck{popdark}}%
  \def\star{\mathbin{\ast}}\def\bigstar{\mathord{\ast}}%
  \def\diamond{\mathbin{\scalebox{0.6}{$\Diamond$}}}%
  \def\bigcup{\mathop{\mathchoice{\raisebox{-0.3em}{\scalebox{1.7}{$\cup$}}}{\scalebox{1.25}{$\cup$}}{\cup}{\cup}}\displaylimits}%
  \def\bigcap{\mathop{\mathchoice{\raisebox{-0.3em}{\scalebox{1.7}{$\cap$}}}{\scalebox{1.25}{$\cap$}}{\cap}{\cap}}\displaylimits}%
  \def\lesseqqgtr{\mathrel{\lessgtr}}\def\bigoplus{\mathop{\oplus}\displaylimits}%
}
\providecommand{\ssi}{si et seulement si} % abréviation fréquente
\AtBeginDocument{\providecommand{\wideparen}{\overparen}} % arc de cercle

%% FIGFIX-BEGIN v5 — correctifs globaux des figures (docs/onbuch-generation/tools/figfix)
\usepackage{adjustbox}\usepackage{etoolbox}\tcbuselibrary{raster}
% toute figure TikZ/pgfplots plus large que la ligne est réduite à la largeur disponible
\BeforeBeginEnvironment{tikzpicture}{\begin{adjustbox}{max width=\linewidth}}
\AfterEndEnvironment{tikzpicture}{\end{adjustbox}}
% légendes pgfplots lisibles par-dessus les courbes
\pgfplotsset{every axis/.append style={legend style={fill=white,fill opacity=0.92,text opacity=1,draw=black!15,inner sep=3pt}}}
% étiquettes d'arêtes (syntaxe edge["..."]) avec fond blanc
\tikzset{every edge quotes/.append style={fill=white,inner sep=1.2pt}}
%% FIGFIX-END
\begin{document}

\pagegarde{Equations et inéquations su 1er degré à une inconnue dans IR}{Mathématiques}{3ème}{Mathématiques – classe de 3ème}{N09}{4 séances (fiche de progression harmonisée)}{Cette leçon outille l'élève pour résoudre les équations de la forme ax + b = 0, les équations qui s'y ramènent (avec parenthèses, fractions, produits nuls) ainsi que les inéquations de la forme ax + b > 0 dans IR. Elle insiste sur la justification des transformations et sur l'interprétation concrète des solutions dans des situations de la vie courante.
\vspace{4pt}
\begin{multicols}{2}\small
\begin{itemize}
\item À la fin de cette leçon, je sais reconnaître une équation du premier degré à une inconnue dans IR
\item résoudre une équation de la forme ax + b = 0 et donner son ensemble de solutions
\item résoudre des équations se ramenant au premier degré (développement, réduction, fractions, produit nul)
\item résoudre une inéquation de la forme ax + b > 0 en appliquant correctement la règle du changement de sens
\item représenter l'ensemble des solutions d'une inéquation sur une droite graduée
\item traduire une situation concrète de la vie courante par une équation ou une inéquation du premier degré
\end{itemize}\end{multicols}}

\begin{sommaire}
\begin{multicols}{2}
\begin{enumerate}
\item Équations de la forme ax + b = 0
\item Équations se ramenant au premier degré
\item Inégalités et règles de transformation
\item Inéquations de la forme ax + b > 0
\item Résolution de problèmes concrets
\item Activité d'intégration
\item Méthodes et exercices résolus
\item Exercices
\item Corrigés des exercices
\item Fiche bilan
\end{enumerate}
\end{multicols}
\end{sommaire}

\begin{objectifs}
\begin{itemize}
\item À la fin de cette leçon, je sais reconnaître une équation du premier degré à une inconnue dans IR
\item résoudre une équation de la forme ax + b = 0 et donner son ensemble de solutions
\item résoudre des équations se ramenant au premier degré (développement, réduction, fractions, produit nul)
\item résoudre une inéquation de la forme ax + b > 0 en appliquant correctement la règle du changement de sens
\item représenter l'ensemble des solutions d'une inéquation sur une droite graduée
\item traduire une situation concrète de la vie courante par une équation ou une inéquation du premier degré
\item rédiger et justifier chaque étape d'une démarche de résolution
\item vérifier qu'une valeur donnée est ou n'est pas solution d'une équation ou d'une inéquation
\end{itemize}
\end{objectifs}

\begin{prerequis}
\begin{itemize}
\item Calcul littéral de 4ème : développer, factoriser, réduire une expression
\item Résolution d'équations simples du type x + a = b et ax = b (4ème)
\item Opérations sur les nombres relatifs et les fractions
\item Comparaison de nombres relatifs et droite graduée
\item Notion d'inégalité et symboles <, >, ≤, ≥
\end{itemize}
\end{prerequis}

\newpage

\coursec{Équations de la forme $ax + b = 0$}

\courssub{Situation d'accroche : Mama Ngo Bell au marché Mokolo}

Imagine \textbf{Mama Ngo Bell} installée sous son parapluie coloré au cœur du marché Mokolo, à Yaoundé. Elle vend de délicieux beignets \emph{puff-puff} à \SI{25}{FCFA} l'unité. Ce matin, elle a dépensé exactement \SI{1500}{FCFA} pour acheter la farine, l'huile, le sucre et la levure. Elle se pose deux questions très concrètes :
\begin{enumerate}
    \item Combien de beignets doit-elle vendre pour \textbf{rembourser exactement ses frais} (ni perte, ni gain) ?
    \item Combien doit-elle vendre \textbf{au minimum} pour commencer à faire du \textbf{bénéfice} ?
\end{enumerate}
Si l'on note $x$ le nombre de beignets vendus, sa recette est $25x$ francs. La première question s'écrit : $25x = 1\,500$, soit $25x - 1\,500 = 0$. La seconde s'écrit : $25x > 1\,500$, soit $25x - 1\,500 > 0$. Nous voilà en présence d'une \textbf{\cle{équation}} et d'une \textbf{\cle{inéquation}} du premier degré à une inconnue. C'est précisément l'objet de cette leçon : apprendre à les résoudre avec rigueur et méthode.

\courssub{Rappels : égalité, inconnue, équation, solution}

Avant d'aller plus loin, posons les bases du vocabulaire. Une \textbf{\cle{égalité}} est une relation qui affirme que deux expressions numériques ont la même valeur (exemple : $3 + 5 = 8$). Une \textbf{\cle{inconnue}} est un nombre que l'on ne connaît pas encore et que l'on représente par une lettre (souvent $x$). Une \textbf{\cle{équation}} est une égalité qui contient au moins une inconnue. \textbf{Résoudre une équation}, c'est trouver toutes les valeurs de l'inconnue qui transforment l'égalité en une vérité numérique. Ces valeurs sont appelées \textbf{\cle{solutions}} de l'équation. L'ensemble de toutes les solutions se note $S$.

\begin{exemplebox}[Équation et solution]
L'égalité $2x + 3 = 9$ est une équation. Si l'on teste $x = 3$, on obtient $2 \times 3 + 3 = 6 + 3 = 9$, l'égalité est vraie : $3$ est une solution. Si l'on teste $x = 2$, on obtient $2 \times 2 + 3 = 7 \neq 9$ : $2$ n'est pas une solution.
\end{exemplebox}

\courssub{Définition d'une équation du premier degré à une inconnue}

Toutes les équations ne se ressemblent pas. Celles qui nous intéressent ici ont une forme bien précise après simplification.

\begin{definition}[Équation du premier degré à une inconnue]
Une \textbf{\cle{équation du premier degré à une inconnue}} est une équation qui, après transformations élémentaires (réduction, factorisation, mise au même dénominateur), peut s'écrire sous la \textbf{forme canonique} :
\[
ax + b = 0
\]
où $a$ et $b$ sont des nombres réels donnés, avec \textbf{$a \neq 0$}, et $x$ est l'inconnue.
\end{definition}

\begin{attention}[Condition $a \neq 0$]
La condition $a \neq 0$ est \textbf{essentielle}. Si $a = 0$, l'équation devient $b = 0$, qui ne contient plus l'inconnue $x$ : ce n'est plus une équation du premier degré à une inconnue (on étudiera ce cas particulier plus loin).
\end{attention}

\courssub{Règles de transformation : la balance à deux plateaux}

Pour résoudre une équation, on la transforme successivement en équations \textbf{\cle{équivalentes}} (qui ont exactement les mêmes solutions) jusqu'à isoler l'inconnue. L'image de la \textbf{balance à deux plateaux} est idéale pour comprendre les règles autorisées : les deux membres de l'égalité sont en équilibre ; toute action identique des deux côtés conserve l'équilibre.

\begin{popfigure}
\begin{tikzpicture}[scale=0.85, every node/.style={font=\small}]
    % Base
    \draw[popdark, thick] (0,0) -- (10,0);
    % Pivot central
    \draw[popdark, thick] (5,0) -- (5,1);
    \draw[popdark, fill=popgold] (5,1) circle (0.3);
    % Bras gauche
    \draw[popdark, thick] (5,1) -- (1,3);
    % Bras droit
    \draw[popdark, thick] (5,1) -- (9,3);
    % Fils
    \draw[popdark, thick] (1,3) -- (1,4.5);
    \draw[popdark, thick] (9,3) -- (9,4.5);
    % Plateaux
    \draw[popblue, thick, fill=popblue!10] (0,4.5) rectangle (2,5);
    \draw[poporange, thick, fill=poporange!10] (8,4.5) rectangle (10,5);
    % Étiquettes plateaux
    \node[popblue] at (1,4.75) {\textbf{Membre gauche}};
    \node[poporange] at (9,4.75) {\textbf{Membre droit}};
    % Actions autorisées
    \node[popgreen, align=center, font=\small] at (11.5,4.2) {Ajouter $c$\\(ou retrancher $c$)};
    \node[poppurple, align=center, font=\small] at (11.5,3) {Multiplier par $k \neq 0$\\(ou diviser par $k \neq 0$)};
    % Flèches vers les plateaux
    \draw[->, popgreen, thick] (10.5,4.2) -- (10,4.2);
    \draw[->, poppurple, thick] (10.5,3) -- (10,3);
    % Poids symboliques
    \node[popmuted, rotate=90] at (1, 3.8) {même poids};
    \node[popmuted, rotate=90] at (9, 3.8) {même poids};
\end{tikzpicture}
\legende{La balance de l'égalité : ajouter/retirer le même poids des deux côtés, ou multiplier/diviser les poids par le même coefficient non nul, conserve l'équilibre.}
\end{popfigure}

\begin{propriete}[Règles de transformation d'une égalité]
Soit l'égalité $A = B$ où $A$ et $B$ sont des expressions contenant l'inconnue $x$, et $c, k \in \mathbb{R}$.
\begin{enumerate}
    \item \textbf{Ajouter (ou retrancher) un même nombre aux deux membres} : \\
    $A = B \;\Longleftrightarrow\; A + c = B + c \;\Longleftrightarrow\; A - c = B - c$.
    \item \textbf{Multiplier (ou diviser) les deux membres par un même nombre non nul} : \\
    $A = B \;\Longleftrightarrow\; kA = kB \quad (k \neq 0) \;\Longleftrightarrow\; \dfrac{A}{k} = \dfrac{B}{k} \quad (k \neq 0)$.
\end{enumerate}
Ces transformations produisent des équations \textbf{équivalentes} : elles ont exactement le même ensemble de solutions.
\end{propriete}

\courssub{Résolution de $ax + b = 0$ (cas $a \neq 0$)}

Appliquons ces règles à la forme canonique $ax + b = 0$ avec $a \neq 0$. L'objectif est d'isoler $x$.

\begin{experience}[Démonstration guidée : isolation de $x$]
Partons de $ax + b = 0$.
\begin{enumerate}
    \item \textbf{Transposer $b$ à droite} (on retranche $b$ des deux membres) : \\
    $ax + b - b = 0 - b \quad\Longrightarrow\quad ax = -b$.
    \item \textbf{Diviser par $a$} (on divise les deux membres par $a$, ce qui est autorisé car $a \neq 0$) : \\
    $\dfrac{ax}{a} = \dfrac{-b}{a} \quad\Longrightarrow\quad x = -\dfrac{b}{a}$.
\end{enumerate}
\emph{Conclusion :} L'équation $ax + b = 0$ (avec $a \neq 0$) admet une \textbf{unique solution} : $x = -\dfrac{b}{a}$.
\end{experience}

\begin{propriete}[Solution de l'équation $ax + b = 0$]
L'équation $ax + b = 0$, avec $a, b \in \mathbb{R}$ et $a \neq 0$, admet une \textbf{unique solution} :
\[
x = -\frac{b}{a}
\]
L'ensemble des solutions se note $S = \left\{-\dfrac{b}{a}\right\}$.
\end{propriete}

\begin{methode}[Résoudre une équation de la forme $ax + b = 0$]
\begin{enumerate}
    \item \textbf{Identifier} les coefficients $a$ (coefficient de $x$) et $b$ (terme constant).
    \item \textbf{Vérifier} que $a \neq 0$ (sinon, cas particulier).
    \item \textbf{Transposer $b$} : écrire $ax = -b$.
    \item \textbf{Diviser par $a$} : écrire $x = -\dfrac{b}{a}$.
    \item \textbf{Conclure} en donnant l'ensemble solution : $S = \left\{-\dfrac{b}{a}\right\}$.
    \item \textbf{Vérifier} (facultatif mais recommandé) en remplaçant $x$ par la solution dans l'équation initiale.
\end{enumerate}
\end{methode}

\begin{exemplebox}[Résolution de $3x + 12 = 0$]
\begin{enumerate}
    \item Identification : $a = 3$, $b = 12$. On a bien $a \neq 0$.
    \item Transposition : $3x = -12$.
    \item Division par $3$ : $x = -\dfrac{12}{3} = -4$.
    \item Conclusion : $S = \{-4\}$.
    \item \textbf{Vérification} : $3 \times (-4) + 12 = -12 + 12 = 0$. L'égalité est vérifiée.
\end{enumerate}
\end{exemplebox}

\begin{attention}[Pièges fréquents]
\begin{itemize}
    \item \textbf{Oublier le signe moins} : La solution est $x = -\dfrac{b}{a}$, pas $\dfrac{b}{a}$. Exemple : $2x - 6 = 0 \Rightarrow 2x = 6 \Rightarrow x = 3$ (et non $-3$). Ici $b = -6$, donc $-\frac{b}{a} = -\frac{-6}{2} = 3$.
    \item \textbf{Diviser par $a$ sans vérifier $a \neq 0$} : Si $a = 0$, la division est interdite et le raisonnement s'effondre.
    \item \textbf{Mal rédiger la conclusion} : On écrit $S = \{-4\}$ (ensemble contenant le nombre $-4$), pas $x = -4$ seul, et certainement pas $S = -4$ (un ensemble n'est pas un nombre).
\end{itemize}
\end{attention}

\courssub{Cas particuliers : que se passe-t-il si $a = 0$ ?}

La définition impose $a \neq 0$. Mais dans un exercice, rien ne garantit que le coefficient devant $x$ soit non nul après simplification. Il faut \textbf{toujours} regarder le coefficient de $x$ avant de diviser.

\begin{propriete}[Cas $a = 0$ dans $ax + b = 0$]
Soit l'équation $ax + b = 0$ avec $a = 0$. Elle se réduit à $b = 0$.
\begin{itemize}
    \item Si $b \neq 0$ (exemple : $0x + 5 = 0 \Leftrightarrow 5 = 0$), l'égalité est \textbf{fausse} quel que soit $x$. \\
    \textbf{Aucune solution} : $S = \varnothing$ (ensemble vide).
    \item Si $b = 0$ (exemple : $0x + 0 = 0 \Leftrightarrow 0 = 0$), l'égalité est \textbf{toujours vraie} quel que soit $x$. \\
    \textbf{Tout réel est solution} : $S = \mathbb{R}$.
\end{itemize}
\end{propriete}

\begin{exoresolu}[Cas particuliers]
Résoudre les équations suivantes :
\begin{enumerate}
    \item $0 \cdot x + 7 = 0$
    \item $(2 - 2)x - 3 = 0$
    \item $5x - 5x + 0 = 0$
\end{enumerate}
\tcblower
\textbf{Solution :}
\begin{enumerate}
    \item $7 = 0$ est faux. $S = \varnothing$.
    \item $0 \cdot x - 3 = 0 \Leftrightarrow -3 = 0$ faux. $S = \varnothing$.
    \item $0 = 0$ vrai. $S = \mathbb{R}$.
\end{enumerate}
\end{exoresolu}

\courssub{Illustration : la solution sur la droite graduée}

La solution d'une équation du premier degré est un nombre unique (si $a \neq 0$). On peut la repérer sur la droite graduée.

\begin{popfigure}
\begin{tikzpicture}[scale=1.2]
    % Axe principal
    \draw[->, thick, popdark] (-6,0) -- (2,0) node[right] {$x$};
    % Graduations
    \foreach \x in {-5,-4,-3,-2,-1,0,1} {
        \draw[thick] (\x, 0.1) -- (\x, -0.1) node[below] {$\x$};
    }
    % Mise en évidence de la solution x = -4
    \draw[popblue, very thick] (-4, -0.3) -- (-4, 0.8);
    \fill[popblue] (-4, 0.5) circle (0.1);
    \node[popblue, above] at (-4, 0.9) {$x = -4$};
    % Flèches indicatives
    \draw[popmuted, <->] (-5, -0.5) -- (-4, -0.5) node[midway, below] {$1$};
    \draw[popmuted, <->] (-4, -0.5) -- (-3, -0.5) node[midway, below] {$1$};
    % Légende
    \node[align=center, popmuted, font=\small] at (-2, -1.5) {Solution de $3x+12=0$ : $S=\{-4\}$};
\end{tikzpicture}
\legende{Repérage de la solution unique $x = -4$ sur la droite réelle.}
\end{popfigure}

\courssub{Le savais-tu ?}

\begin{savaistu}[Des scribes babyloniens aux algorithmes modernes]
Les équations du premier degré ne datent pas d'hier ! Il y a près de \textbf{4 000 ans}, en Mésopotamie (l'actuel Irak), les scribes babyloniens résolvaient déjà des problèmes du type « J'ai ajouté le côté de mon carré à sa surface, le total fait 0;45 » (en notation sexagésimale). Ils utilisaient des méthodes algorithmiques sur tablettes d'argile, ancêtres de notre résolution $x = -b/a$. Au Cameroun aujourd'hui, que ce soit pour calculer le point mort d'une petite entreprise au marché Mokolo, dimensionner un câble électrique pour CAMTEL, ou programmer un virement Mobile Money (MTN MoMo / Orange Money), on utilise exactement la même logique mathématique : \emph{modéliser par une équation, résoudre, interpréter}.
\end{savaistu}

\courssub{Exercice résolu complet : de l'énoncé à la conclusion}

\begin{exoresolu}[Problème de gestion : le seuil de rentabilité]
Mama Ngo Bell (notre vendeuse de beignets) a amélioré sa recette. Ses frais fixes (location du stand, gaz) s'élèvent à \SI{2000}{FCFA} par jour. Chaque beignet lui coûte \SI{15}{FCFA} à fabriquer (farine, huile, sucre) et elle le revend \SI{50}{FCFA}.
\begin{enumerate}
    \item Exprimer son bénéfice $B$ (en FCFA) en fonction du nombre $x$ de beignets vendus.
    \item Déterminer le nombre de beignets à vendre pour que le bénéfice soit nul (seuil de rentabilité).
    \item Vérifier le résultat.
\end{enumerate}
\tcblower
\textbf{Solution détaillée :}
\begin{enumerate}
    \item \textbf{Modélisation} : Recette = $50x$. Coût total = Frais fixes + Coût variable = $2\,000 + 15x$.
    Bénéfice $B = \text{Recette} - \text{Coût total} = 50x - (2\,000 + 15x) = 35x - 2\,000$.
    \item \textbf{Seuil de rentabilité} : $B = 0 \Leftrightarrow 35x - 2\,000 = 0$.
    C'est une équation de la forme $ax + b = 0$ avec $a = 35 \neq 0$ et $b = -2\,000$.
    \begin{align*}
        35x - 2\,000 &= 0 \\
        35x &= 2\,000 \quad \text{(transposition de } -2\,000\text{)} \\
        x &= \frac{2\,000}{35} \quad \text{(division par } 35\text{)} \\
        x &= \frac{400}{7} \approx 57,14
    \end{align*}
    Comme on ne peut pas vendre une fraction de beignet, Mama Ngo Bell doit vendre \textbf{au moins 58 beignets} pour ne pas perdre d'argent (au 57ème, elle perd encore un peu).
    \item \textbf{Vérification} :
    \begin{itemize}
        \item Pour $x = 57$ : $B = 35 \times 57 - 2\,000 = 1\,995 - 2\,000 = -5$ FCFA (perte de 5 FCFA).
        \item Pour $x = 58$ : $B = 35 \times 58 - 2\,000 = 2\,030 - 2\,000 = 30$ FCFA (bénéfice de 30 FCFA).
    \end{itemize}
    Le seuil mathématique exact est $x = \frac{400}{7}$, le seuil pratique (entier) est $58$.
\end{enumerate}
\end{exoresolu}

\begin{aretenir}[Résumé de la section]
\begin{itemize}
    \item Une équation du 1er degré à une inconnue s'écrit $ax + b = 0$ avec $a \neq 0$.
    \item Règles d'équivalence : on peut ajouter/retirer le même nombre des deux côtés, ou multiplier/diviser par un même nombre \textbf{non nul}.
    \item Résolution : $ax + b = 0 \;\Rightarrow\; ax = -b \;\Rightarrow\; x = -\dfrac{b}{a}$.
    \item Ensemble solution : $S = \left\{-\dfrac{b}{a}\right\}$.
    \item Cas $a = 0$ : si $b \neq 0 \Rightarrow S = \varnothing$ ; si $b = 0 \Rightarrow S = \mathbb{R}$.
    \item Toujours vérifier la solution par substitution dans l'équation initiale.
\end{itemize}
\end{aretenir}

\coursec{Équations se ramenant au premier degré}

Dans la section précédente, nous avons maîtrisé la résolution des équations de la forme canonique $ax + b = 0$ avec $a \neq 0$. Or, dans la vie courante comme dans les problèmes d'examen, les équations ne se présentent que rarement sous cette forme toute faite. Elles contiennent des parenthèses, l'inconnue apparaît des deux côtés de l'égalité, ou bien elles font intervenir des fractions, voire des produits de facteurs. La stratégie universelle en mathématiques, et particulièrement en algèbre, consiste à \textbf{transformer} une situation complexe en une situation simple que l'on sait déjà résoudre. C'est exactement ce que nous allons faire ici : transformer toute équation « compliquée » en une équation du type $ax + b = 0$, la résoudre, puis vérifier que la solution trouvée convient bien à l'équation de départ.

\courssub{Équations avec parenthèses : développer pour simplifier}

La première difficulté fréquente est la présence de parenthèses. Rappelons la propriété fondamentale de la distributivité : pour tous nombres réels $k$, $a$, $b$, on a $k(a + b) = ka + kb$ et $k(a - b) = ka - kb$. Cette propriété permet de « faire sauter » les parenthèses en multipliant le terme extérieur par \textbf{chacun} des termes intérieurs.

\begin{methode}[Résoudre une équation avec parenthèses]
\leavevmode
\begin{enumerate}
    \item \textbf{Développer} chaque membre en appliquant la distributivité (attention aux signes \og moins \fg{} devant une parenthèse !).
    \item \textbf{Réduire} chaque membre en regroupant les termes semblables (termes en $x$ d'un côté, constantes de l'autre).
    \item \textbf{Ramener} l'équation à la forme $ax + b = 0$ par transposition.
    \item \textbf{Résoudre} et \textbf{vérifier} la solution dans l'équation \textbf{initiale} (celle avec les parenthèses).
\end{enumerate}
\end{methode}

\begin{exoresolu}[Équation avec parenthèses : $5(x - 2) = 3x + 4$]
\emph{Énoncé.} Résoudre dans $\mathbb{R}$ l'équation $5(x - 2) = 3x + 4$.

\emph{Solution.}
\begin{align*}
    5(x - 2) &= 3x + 4 & \text{(Équation de départ)} \\
    5 \times x - 5 \times 2 &= 3x + 4 & \text{(Distributivité : $5(x-2)=5x-10$)} \\
    5x - 10 &= 3x + 4 & \text{(Membre de gauche réduit)} \\
    5x - 3x &= 4 + 10 & \text{(Transposition : $5x-3x$ à gauche, $-10$ à droite en changeant de signe)} \\
    2x &= 14 & \text{(Réduction)} \\
    x &= \frac{14}{2} & \text{(Division par le coefficient de $x$)} \\
    x &= 7
\end{align*}

\emph{Vérification :} Remplaçons $x$ par $7$ dans l'équation \textbf{initiale}.
Membre de gauche : $5(7 - 2) = 5 \times 5 = 25$.
Membre de droite : $3 \times 7 + 4 = 21 + 4 = 25$.
Les deux membres sont égaux, donc $7$ est bien solution.
\emph{Conclusion :} L'ensemble solution est $S = \{7\}$.
\tcblower
\textbf{Analyse pédagogique :} L'étape cruciale est le développement $5(x-2) = 5x - 10$. Une erreur très classique est d'écrire $5x - 2$ (oubli de multiplier le $-2$ par $5$). Vérifiez toujours en remplaçant dans l'équation \emph{non développée}.
\end{exoresolu}

\begin{attention}[Piège mortel : la distributivité oubliée]
\leavevmode
L'erreur n°1 au BEPC sur ce chapitre est d'écrire :
\[ 2(x + 3) = 2x + 3 \quad \text{(FAUX !)} \]
Le $2$ doit multiplier \textbf{les deux} termes de la parenthèse :
\[ 2(x + 3) = 2 \times x + 2 \times 3 = 2x + 6 \quad \text{(VRAI)} \]
De même, attention au signe moins devant une parenthèse : $-(x - 4) = -x + 4$ (le moins change les signes de \emph{tous} les termes intérieurs).
\end{attention}

\courssub{Équations avec l'inconnue dans les deux membres : regrouper}

Lorsque l'inconnue $x$ apparaît des deux côtés du signe $=$, le but est de la \textbf{regrouper dans un seul membre} (généralement celui où son coefficient sera positif pour éviter les erreurs de signe). On utilise les transformations qui conservent l'égalité : ajouter ou soustraire le même terme des deux côtés.

\begin{exemplebox}[Regroupement des termes en $x$ : $4x - 7 = 2x + 5$]
\leavevmode
On veut réunir les $x$ à gauche (coefficient $4 > 2$).
\begin{align*}
    4x - 7 &= 2x + 5 \\
    4x - 2x - 7 &= 5 & \text{(On soustrait $2x$ des deux côtés)} \\
    2x - 7 &= 5 \\
    2x &= 5 + 7 & \text{(On ajoute $7$ des deux côtés)} \\
    2x &= 12 \\
    x &= 6
\end{align*}
\emph{Vérification :} $4(6)-7 = 24-7=17$ et $2(6)+5=12+5=17$. \quad $S = \{6\}$.
\end{exemplebox}

\courssub{Équations fractionnaires : éliminer les dénominateurs}

Une équation fractionnaire contient l'inconnue au numérateur avec des dénominateurs constants (ex: $\frac{x+1}{2} = \frac{3x-2}{5}$) ou simplement des fractions. La technique standard consiste à multiplier les deux membres par le \textbf{PPCM} (Plus Petit Commun Multiple) des dénominateurs pour obtenir une équation sans fractions, donc du premier degré.

\begin{propriete}[Élimination des dénominateurs]
\leavevmode
Soit une équation où les deux membres sont des expressions rationnelles à dénominateurs constants non nuls. Si l'on multiplie les deux membres par un même nombre non nul (notamment le PPCM des dénominateurs), on obtient une équation \textbf{équivalente} (qui a les mêmes solutions).
\end{propriete}

\begin{methode}[Résoudre une équation fractionnaire (dénominateurs constants)]
\leavevmode
\begin{enumerate}
    \item Identifier les dénominateurs et calculer leur \textbf{PPCM}.
    \item \textbf{Multiplier les deux membres} de l'équation par ce PPCM.
    \item Simplifier chaque fraction (le dénominateur doit disparaître).
    \item Résoudre l'équation du premier degré obtenue.
    \item \textbf{Vérifier} la solution dans l'équation initiale.
\end{enumerate}
\end{methode}

\begin{exoresolu}[Équation fractionnaire : $\frac{x+1}{3} - \frac{2x-5}{2} = 1$]
\emph{Énoncé.} Résoudre dans $\mathbb{R}$ : $\frac{x+1}{3} - \frac{2x-5}{2} = 1$.

\emph{Solution.}
Les dénominateurs sont $3$ et $2$. Leur PPCM est $6$.
On multiplie les deux membres par $6$ :
\[ 6 \times \left( \frac{x+1}{3} - \frac{2x-5}{2} \right) = 6 \times 1 \]
On distribue le $6$ à gauche :
\[ 6 \times \frac{x+1}{3} - 6 \times \frac{2x-5}{2} = 6 \]
On simplifie : $\frac{6}{3}=2$ et $\frac{6}{2}=3$.
\[ 2(x+1) - 3(2x-5) = 6 \]
On développe (attention au moins devant $3$) :
\[ 2x + 2 - 6x + 15 = 6 \]
On regroupe :
\[ -4x + 17 = 6 \]
\[ -4x = 6 - 17 \]
\[ -4x = -11 \]
\[ x = \frac{-11}{-4} = \frac{11}{4} = 2,75 \]

\emph{Vérification :}
Membre de gauche : $\frac{2,75+1}{3} - \frac{2(2,75)-5}{2} = \frac{3,75}{3} - \frac{5,5-5}{2} = 1,25 - \frac{0,5}{2} = 1,25 - 0,25 = 1$.
Membre de droite : $1$. \quad $S = \{\frac{11}{4}\}$.
\tcblower
\textbf{Astuce :} Écrire la multiplication par le PPCM \emph{au-dessus} de la barre de fraction aide à visualiser les simplifications : $\frac{\cancel{6}^2(x+1)}{\cancel{3}} - \frac{\cancel{6}^3(2x-5)}{\cancel{2}} = 6$.
\end{exoresolu}

\begin{attention}[Erreur fréquente : oublier de multiplier \textbf{tous} les termes]
\leavevmode
Si l'équation est $\frac{x}{2} + 3 = 5$, multiplier par $2$ donne :
\[ 2 \times \left( \frac{x}{2} + 3 \right) = 2 \times 5 \implies x + 6 = 10 \]
\emph{Pas :} $x + 3 = 10$ (le $3$ a été oublié). Le PPCM doit multiplier \textbf{chaque terme} de l'équation, pas seulement les fractions.
\end{attention}

\courssub{Équations-produits : la règle du produit nul}

Nous rencontrons maintenant un type d'équation radicalement différent : le premier membre est un \textbf{produit} de deux facteurs contenant l'inconnue, et le second membre est $0$. Par exemple : $(2x - 6)(x + 1) = 0$.

\begin{propriete}[Règle du produit nul (Admise en 3ème)]
\leavevmode
Soit $A$ et $B$ deux expressions littérales (ou nombres). Le produit $A \times B$ est nul \textbf{si et seulement si} l'un au moins des facteurs est nul.
\[ A \times B = 0 \quad \Longleftrightarrow \quad A = 0 \quad \text{ou} \quad B = 0 \]
\textit{Justification intuitive :} Si tu multiplies deux nombres et que le résultat est $0$, c'est que \emph{forcément} l'un des deux était $0$. Il n'y a pas d'autre possibilité dans $\mathbb{R}$.
\end{propriete}

\begin{methode}[Résoudre une équation du type $(ax+b)(cx+d)=0$]
\leavevmode
\begin{enumerate}
    \item Appliquer la règle du produit nul : écrire les deux équations séparées $ax+b=0$ \textbf{et} $cx+d=0$.
    \item Résoudre \textbf{chacune} de ces deux équations du premier degré.
    \item \textbf{Réunir} les solutions obtenues dans un ensemble $S$.
    \item Vérifier (facultatif mais recommandé) chaque solution dans l'équation produit initiale.
\end{enumerate}
\end{methode}

\begin{exemplebox}[Équation-produit : $(2x - 6)(x + 1) = 0$]
\leavevmode
\emph{Énoncé.} Résoudre dans $\mathbb{R}$ : $(2x - 6)(x + 1) = 0$.

\emph{Solution.}
D'après la règle du produit nul :
\[ (2x - 6)(x + 1) = 0 \quad \Longleftrightarrow \quad \begin{cases} 2x - 6 = 0 \\ \text{ou} \\ x + 1 = 0 \end{cases} \]

\emph{Première équation :} $2x - 6 = 0 \implies 2x = 6 \implies x = 3$.
\emph{Deuxième équation :} $x + 1 = 0 \implies x = -1$.

Les deux valeurs conviennent. L'ensemble solution est $S = \{-1 ; 3\}$.
\tcblower
\textbf{Remarque importante :} On écrit \og ou \fg{} entre les deux équations car $x$ ne peut pas être \emph{en même temps} $3$ et $-1$. C'est une \textbf{réunion} de solutions, pas une intersection.
\end{exemplebox}

\begin{attention}[Piège : diviser par un facteur au lieu d'appliquer la règle du produit nul]
\leavevmode
Face à $(x-3)(x+2)=0$, l'erreur fatale est de diviser les deux membres par $(x-3)$ :
\[ \frac{(x-3)(x+2)}{x-3} = \frac{0}{x-3} \implies x+2=0 \implies x=-2 \]
\emph{Vous venez de perdre la solution $x=3$ !} Car diviser par $(x-3)$ suppose que $x-3 \neq 0$, ce qui exclut $x=3$ \emph{a priori}.
\textbf{Jamais de division par une expression contenant l'inconnue.} Utilisez toujours la règle du produit nul.
\end{attention}

\begin{savaistu}[La règle du produit nul : la clé du second degré]
\leavevmode
Tu te demandes peut-être pourquoi on t'apprend la règle du produit nul $(A \times B = 0 \iff A=0 \text{ ou } B=0)$ en 3ème alors qu'on ne résout pas encore d'équations du second degré ? C'est parce que \textbf{toute} équation du second degré $ax^2+bx+c=0$ (avec $a \neq 0$) se résout au niveau supérieur \emph{exclusivement} en factorisant le trinôme pour le transformer en produit nul : $(x-x_1)(x-x_2)=0$. Maîtriser parfaitement cette règle aujourd'hui, c'est t'assurer une base solide pour l'algèbre des classes supérieures ! Au Cameroun, comme partout, les élèves qui rencontrent des difficultés par la suite sont souvent ceux qui n'ont pas automatisé : \og Produit nul $\rightarrow$ Facteurs nuls \fg{}.
\end{savaistu}

\courssub{Démarche générale de résolution : l'organigramme de synthèse}

Pour ne jamais te perdre face à une équation inconnue, suis cet organigramme. Il résume \textbf{toute} la logique de cette section : transformer pour revenir au connu ($ax+b=0$), sauf cas particulier du produit nul.

\begin{popfigure}
\begin{tikzpicture}[
    node distance=1.2cm and 1.5cm,
    startstop/.style={rectangle, rounded corners, minimum width=3cm, minimum height=1cm, text centered, draw=popblue, fill=popblueL, thick},
    process/.style={rectangle, minimum width=3.5cm, minimum height=0.9cm, text centered, draw=popdark, fill=popcream, thick},
    decision/.style={diamond, aspect=2, minimum width=3cm, minimum height=1cm, text centered, draw=poporange, fill=poporangeL, thick},
    io/.style={trapezium, trapezium left angle=70, trapezium right angle=110, minimum width=3cm, minimum height=1cm, text centered, draw=popgreen, fill=popgreenL, thick},
    arrow/.style={->, >=Stealth, thick, popdark},
    yesno/.style={midway, fill=white, inner sep=2pt, font=\small}
]
% Nœuds
\node (start) [startstop] {Équation donnée};
\node (dec1) [decision, below of=start] {Y a-t-il des parenthèses ?};
\node (proc1) [process, below of=dec1] {Développer (distributivité)};
\node (dec2) [decision, below of=proc1] {Y a-t-il des dénominateurs ?};
\node (proc2) [process, below of=dec2] {Multiplier par le PPCM des dénominateurs};
\node (dec3) [decision, below of=proc2] {Le 1er membre est-il un produit \textbf{égal à 0} ?};
\node (proc3a) [process, left of=dec3, xshift=-3.5cm, align=center]{Règle du produit nul :\\ $A \times B = 0 \Rightarrow A=0$ ou $B=0$\\ Résoudre chaque équation du 1er degré};
\node (proc3b) [process, right of=dec3, xshift=3.5cm, align=center]{Regrouper les $x$ à gauche, constantes à droite\\ Réduire à la forme $ax + b = 0$};
\node (proc4) [process, below of=dec3, yshift=-2.5cm] {Résoudre $ax+b=0 \Rightarrow x = -\frac{b}{a}$};
\node (verif) [io, below of=proc4] {Vérifier dans l'équation \textbf{initiale}};
\node (end) [startstop, below of=verif] {Écrire l'ensemble solution $S$};

% Flèches
\draw [arrow] (start) -- (dec1);
\draw [arrow] (dec1) -- node[yesno, left] {Oui} (proc1);
\draw [arrow] (dec1) -- node[yesno, right] {Non} (dec2);
\draw [arrow] (proc1) -- (dec2);
\draw [arrow] (dec2) -- node[yesno, left] {Oui} (proc2);
\draw [arrow] (dec2) -- node[yesno, right] {Non} (dec3);
\draw [arrow] (proc2) -- (dec3);
\draw [arrow] (dec3) -- node[yesno, above] {Oui} (proc3a);
\draw [arrow] (dec3) -- node[yesno, above] {Non} (proc3b);
\draw [arrow] (proc3a) |- (verif);
\draw [arrow] (proc3b) -- (proc4);
\draw [arrow] (proc4) -- (verif);
\draw [arrow] (verif) -- (end);
\end{tikzpicture}
\legende{Organigramme de la démarche unifiée de résolution d'une équation se ramenant au premier degré.}
\end{popfigure}

\begin{aretenir}[Démarche type pour toute équation du 1er degré]
\leavevmode
Pour résoudre n'importe quelle équation du premier degré à une inconnue :
\begin{enumerate}
    \item \textbf{Nettoyer} : Développer les parenthèses, éliminer les dénominateurs (PPCM).
    \item \textbf{Cas particulier} : Si l'équation est un \textbf{produit nul} $(A)(B)=0$, appliquer la règle du produit nul (deux équations simples).
    \item \textbf{Cas général} : Regrouper les termes en $x$ dans un membre, les nombres dans l'autre. Réduire à $ax+b=0$.
    \item \textbf{Résoudre} : $x = -\frac{b}{a}$ (si $a \neq 0$).
    \item \textbf{Vérifier} : Remplacer $x$ dans l'équation \textbf{de départ} (pas dans la forme réduite).
    \item \textbf{Conclure} : Écrire $S = \{ \dots \}$ ou $S = \varnothing$ (si $0x = b$ avec $b \neq 0$) ou $S = \mathbb{R}$ (si $0x = 0$).
\end{enumerate}
\end{aretenir}

\begin{exoresolu}[Synthèse : $\frac{3(x-2)}{2} - \frac{x+1}{3} = x - 4$]
\emph{Énoncé.} Résoudre dans $\mathbb{R}$ : $\frac{3(x-2)}{2} - \frac{x+1}{3} = x - 4$.

\emph{Solution.}
\textbf{1. Dénominateurs :} $2$ et $3$. PPCM $= 6$.
Multiplier par $6$ :
\[ 6 \times \frac{3(x-2)}{2} - 6 \times \frac{x+1}{3} = 6(x - 4) \]
Simplifier :
\[ 3 \times 3(x-2) - 2(x+1) = 6x - 24 \]
\[ 9(x-2) - 2(x+1) = 6x - 24 \]

\textbf{2. Développer :}
\[ 9x - 18 - 2x - 2 = 6x - 24 \]
\textbf{3. Réduire membre de gauche :}
\[ 7x - 20 = 6x - 24 \]
\textbf{4. Regrouper $x$ à gauche :}
\[ 7x - 6x = -24 + 20 \]
\[ x = -4 \]

\textbf{5. Vérification (dans l'équation initiale) :}
Membre de gauche : $\frac{3(-4-2)}{2} - \frac{-4+1}{3} = \frac{3(-6)}{2} - \frac{-3}{3} = \frac{-18}{2} + 1 = -9 + 1 = -8$.
Membre de droite : $-4 - 4 = -8$.
Égalité vérifiée.
\emph{Conclusion :} $S = \{-4\}$.
\tcblower
\textbf{Commentaire :} Cet exercice combine \textbf{toutes} les difficultés : parenthèses, fractions, inconnue des deux côtés. La rigueur dans l'ordre des opérations (PPCM $\rightarrow$ Distributivité $\rightarrow$ Regroupement) est la seule garantie contre l'erreur.
\end{exoresolu}

\begin{exercice}[Entraînement ciblé]
\leavevmode
Résoudre dans $\mathbb{R}$ les équations suivantes en détaillant les étapes (développement, regroupement, résolution, vérification, conclusion).
\begin{enumerate}
    \item $4(2x - 1) - 3(x + 5) = 2$
    \item $\frac{x}{4} + \frac{x-2}{3} = 1$
    \item $(x - 4)(2x + 3) = 0$
    \item $5 - 2(3x - 1) = 3(x + 2) - 4x$
    \item $\frac{2x+1}{5} - 1 = \frac{x-3}{2}$
\end{enumerate}
\end{exercice}

\begin{corrige}[Exercice n°1]
\leavevmode
\begin{enumerate}
    \item $8x - 4 - 3x - 15 = 2 \implies 5x - 19 = 2 \implies 5x = 21 \implies x = 4,2$. \quad $S = \{4,2\}$.
    \item PPCM $12$ : $3x + 4(x-2) = 12 \implies 3x + 4x - 8 = 12 \implies 7x = 20 \implies x = \frac{20}{7}$. \quad $S = \{\frac{20}{7}\}$.
    \item Produit nul : $x-4=0$ ou $2x+3=0 \implies x=4$ ou $x=-\frac{3}{2}$. \quad $S = \{-\frac{3}{2} ; 4\}$.
    \item $5 - 6x + 2 = 3x + 6 - 4x \implies 7 - 6x = -x + 6 \implies 7 - 6 = 6x - x \implies 1 = 5x \implies x = 0,2$. \quad $S = \{0,2\}$.
    \item PPCM $10$ : $2(2x+1) - 10 = 5(x-3) \implies 4x+2-10 = 5x-15 \implies 4x-8 = 5x-15 \implies -x = -7 \implies x = 7$. \quad $S = \{7\}$.
\end{enumerate}
\end{corrige}

\coursec{Inégalités et règles de transformation}

\courssub{Des équations aux inéquations : une transition naturelle}

Dans la section précédente, nous avons appris à résoudre des \emph{équations} du premier degré : nous cherchions \textbf{une unique valeur} de $x$ qui rendait l'égalité vraie. Par exemple, $2x + 5 = 11$ admet pour seule solution $x = 3$.

Mais la vie courante, et les mathématiques, nous confrontent souvent à des situations où l'on ne cherche pas une valeur exacte, mais \textbf{toutes les valeurs} qui satisfont une condition de comparaison. Un commerçant du marché Mokolo veut savoir à partir de quel prix de vente il fait \emph{au moins} 5000 FCFA de bénéfice. Un élève veut connaître les notes qu'il doit obtenir au prochain devoir pour que sa moyenne \emph{soit supérieure ou égale} à 12. C'est là qu'interviennent les \textbf{inéquations}.

Une inéquation à une inconnue $x$ s'écrit sous la forme $ax + b > 0$, $ax + b < 0$, $ax + b \ge 0$ ou $ax + b \le 0$. La résoudre, ce n'est pas trouver \emph{un} nombre, mais déterminer \textbf{l'ensemble de tous les nombres réels} qui la vérifient. Cet ensemble est souvent un \emph{intervalle} (une portion de la droite graduée), ce qui signifie qu'il y a une \textbf{infinité de solutions}.

\courssub{Rappel : symboles et lecture d'une inégalité}

Avant de manipuler les inéquations, assurons-nous de parler le même langage. Soient $a$ et $b$ deux nombres réels.

\begin{definition}[Symboles d'ordre]
On note :
\begin{itemize}
    \item $a < b$ : « $a$ est \cle{strictement inférieur} à $b$ » (ou « $b$ est strictement supérieur à $a$ »).
    \item $a > b$ : « $a$ est \cle{strictement supérieur} à $b$ ».
    \item $a \le b$ : « $a$ est \cle{inférieur ou égal} à $b$ ».
    \item $a \ge b$ : « $a$ est \cle{supérieur ou égal} à $b$ ».
\end{itemize}
\end{definition}

\begin{aretenir}[Interprétation géométrique sur la droite graduée]
Sur une droite graduée orientée de la gauche vers la droite :
\begin{itemize}
    \item $a < b$ signifie que le point d'abscisse $a$ se trouve \textbf{à gauche} du point d'abscisse $b$.
    \item $a > b$ signifie que $a$ se trouve \textbf{à droite} de $b$.
\end{itemize}
L'ordre « naturel » des nombres correspond à leur position de la gauche vers la droite.
\end{aretenir}

\begin{exemplebox}[Lecture et positionnement]
\begin{itemize}
    \item $-4 < 2$ : sur la droite, $-4$ est bien à gauche de $2$.
    \item $5 \ge 5$ : c'est vrai (égalité).
    \item $0 > -3,5$ : $0$ est à droite de $-3,5$.
\end{itemize}
\end{exemplebox}

\courssub{Les trois règles d'or de la transformation des inégalités}

Pour résoudre une inéquation comme $3x - 7 < 2$, nous allons la transformer étape par étape en inéquations \emph{équivalentes} (qui ont exactement le même \cle{ensemble des solutions}) jusqu'à isoler $x$. Trois règles fondamentales régissent ces transformations. Elles découlent toutes de la définition : $a < b \iff b - a > 0$ (la différence est positive).

\begin{propriete}[Règles de transformation des inégalités]
Soient $a$, $b$ et $c$ trois nombres réels.
\begin{enumerate}
    \item \textbf{Ajouter (ou retrancher) un même nombre aux deux membres} ne change pas le sens de l'inégalité.
    \[
    a < b \iff a + c < b + c \quad \text{et} \quad a - c < b - c
    \]
    \item \textbf{Multiplier (ou diviser) les deux membres par un nombre \textbf{strictement positif}} ne change pas le sens de l'inégalité.
    \[
    \text{Si } c > 0 : \quad a < b \iff ac < bc \quad \text{et} \quad \frac{a}{c} < \frac{b}{c}
    \]
    \item \textbf{Multiplier (ou diviser) les deux membres par un nombre \textbf{strictement négatif}} \textbf{change le sens} de l'inégalité ($<$ devient $>$, $\le$ devient $\ge$, etc.).
    \[
    \text{Si } c < 0 : \quad a < b \iff ac > bc \quad \text{et} \quad \frac{a}{c} > \frac{b}{c}
    \]
\end{enumerate}
Ces règles restent valides si l'on remplace $<$ par $>$, $\le$ ou $\ge$ (en adaptant le sens final pour la règle 3).
\end{propriete}

\begin{experience}[Démonstration guidée des règles à partir de la définition]
Rappelons la définition : $a < b$ signifie exactement que $b - a > 0$ (la différence est un nombre strictement positif).

\textbf{1. Règle de l'addition (Règle 1).}
On veut comparer $a+c$ et $b+c$. Calculons leur différence :
\[
(b+c) - (a+c) = b + c - a - c = b - a.
\]
Or, par hypothèse, $b - a > 0$. Donc $(b+c) - (a+c) > 0$, ce qui signifie $a+c < b+c$. \textbf{Le sens est conservé.}

\textbf{2. Règle de la multiplication par un positif (Règle 2).}
Soit $c > 0$. Comparons $ac$ et $bc$ :
\[
bc - ac = c(b - a).
\]
Produit de deux nombres strictement positifs ($c>0$ et $b-a>0$) : le résultat est $> 0$.
Donc $bc - ac > 0$, c'est-à-dire $ac < bc$. \textbf{Le sens est conservé.}

\textbf{3. Règle de la multiplication par un négatif (Règle 3).}
Soit $c < 0$. Comparons $ac$ et $bc$ :
\[
bc - ac = c(b - a).
\]
Produit d'un nombre strictement négatif ($c<0$) et d'un nombre strictement positif ($b-a>0$) : le résultat est $< 0$.
Donc $bc - ac < 0$, c'est-à-dire $bc < ac$, ou équivalentement $ac > bc$. \textbf{Le sens est inversé.}
\end{experience}

\begin{popfigure}
\begin{tikzpicture}[scale=1.2, >=Stealth]
    % Droite graduée originale
    \draw[thick, <->] (-4, 2) -- (5, 2);
    \foreach \x in {-3,-2,-1,0,1,2,3,4}
        \draw (\x, 1.9) -- (\x, 2.1) node[below=2pt] {\small $\x$};
    
    % Points -2 et 3
    \fill[popblue] (-2, 2) circle (3pt) node[above=4pt, popblue] {$A(-2)$};
    \fill[poporange] (3, 2) circle (3pt) node[above=4pt, poporange] {$B(3)$};
    \draw[popblue, thick, <->] (-2, 2.4) -- (3, 2.4) node[midway, above] {$5$ unités};
    \node at (0.5, 2.8) {\textbf{Avant :} $-2 < 3$};

    % Flèche de transformation
    \draw[poppurple, thick, ->] (2.5, 1.5) -- (2.5, 0.5) node[midway, right] {\small $\times (-1)$};

    % Droite graduée transformée
    \draw[thick, <->] (-4, -0.5) -- (5, -0.5);
    \foreach \x in {-3,-2,-1,0,1,2,3,4}
        \draw (\x, -0.6) -- (\x, -0.4) node[below=2pt] {\small $\x$};

    % Points 2 et -3 (images par x -> -x)
    \fill[popblue] (2, -0.5) circle (3pt) node[below=4pt, popblue] {$A'(2)$};
    \fill[poporange] (-3, -0.5) circle (3pt) node[below=4pt, poporange] {$B'(-3)$};
    \draw[poporange, thick, <->] (-3, -0.9) -- (2, -0.9) node[midway, below] {$5$ unités};
    \node at (0.5, -1.3) {\textbf{Après :} $2 > -3$ (sens inversé)};
\end{tikzpicture}
\legende{Multiplier par $-1$ fait faire un demi-tour sur la droite graduée : l'ordre s'inverse.}
\end{popfigure}

\courssub{Exemple chiffré : application en chaîne des règles}

Appliquons les trois règles successivement sur une inégalité simple pour bien voir l'effet de chacune.

\begin{exemplebox}[De $-4 < 6$ à la transformation complète]
Partons de l'inégalité vraie : $\mathbf{-4 < 6}$.

\begin{enumerate}
    \item \textbf{Ajoutons $2$ aux deux membres (Règle 1) :}
    $-4 + 2 < 6 + 2 \iff \mathbf{-2 < 8}$. \quad $\checkmark$ Vrai. Le sens ne change pas.
    
    \item \textbf{Multiplions par $3$ (positif) (Règle 2) :}
    $3 \times (-2) < 3 \times 8 \iff \mathbf{-6 < 24}$. \quad $\checkmark$ Vrai. Le sens ne change pas.
    
    \item \textbf{Multiplions par $-2$ (négatif) (Règle 3) :}
    $(-2) \times (-6) \mathbf{>} (-2) \times 24 \iff \mathbf{12 > -48}$. \quad $\checkmark$ Vrai. \textbf{Le sens a changé} ($<$ est devenu $>$).
\end{enumerate}
\textbf{Bilan :} $-4 < 6 \implies -2 < 8 \implies -6 < 24 \implies 12 > -48$.
\end{exemplebox}

\begin{attention}[Le piège n°1 du BEPC : l'oubli du changement de sens]
C'est \textbf{l'erreur la plus fréquente} et la plus sévèrement sanctionnée.
\begin{center}
    \fcolorbox{poporange}{popcream}{\parbox{0.9\linewidth}{\centering
    \textbf{Quand tu divises (ou multiplies) par un nombre négatif, tu DOIS inverser le signe de l'inégalité.}
    }}
\end{center}
\textbf{Exemple piège :} Résoudre $-2x < 6$.
\begin{itemize}
    \item \textcolor{poporange}{\textbf{Mauvaise résolution :}} $x < \frac{6}{-2} \iff x < -3$. \quad $\rightarrow$ \textbf{FAUX} (sens non inversé).
    \item \textcolor{popgreen}{\textbf{Bonne résolution :}} $x > \frac{6}{-2} \iff x > -3$. \quad $\rightarrow$ \textbf{VRAI}.
\end{itemize}
\textbf{Astuce de pro :} Pour éviter ce piège, \textbf{ne divise jamais par un négatif directement}. Ramène le terme en $x$ du côté où il devient positif en ajoutant/soustrayant.
\[
-2x < 6 \iff 0 < 2x + 6 \iff -6 < 2x \iff -3 < x \iff x > -3.
\]
Ainsi, tu divises toujours par un positif ($+2$ ici) et le sens ne change jamais !
\end{attention}

\courssub{Notion de solution d'une inéquation : une infinité de valeurs}

Contrairement à une équation du type $ax+b=0$ qui a \textbf{une unique solution} (sauf cas particuliers), une inéquation du type $ax+b < 0$ admet \textbf{une infinité de solutions}.

\begin{definition}[Solution d'une inéquation]
Une valeur de $x$ est une \cle{solution} de l'inéquation si, en la remplaçant dans l'inéquation, on obtient une inégalité \textbf{vraie}.
L'\textbf{ensemble des solutions} (noté $S$) est l'ensemble de tous les réels qui vérifient l'inéquation.
\end{definition}

\begin{exemplebox}[Visualiser l'ensemble des solutions]
Prenons l'inéquation $x > -3$.
\begin{itemize}
    \item $-2$ est solution car $-2 > -3$ (vrai).
    \item $0$ est solution car $0 > -3$ (vrai).
    \item $1000$ est solution.
    \item $-3$ \textbf{n'est pas} solution (inégalité stricte).
    \item $-5$ n'est pas solution car $-5 > -3$ est faux.
\end{itemize}
L'ensemble des solutions est l'intervalle \textbf{ouvert} $]-3 ; +\infty[$.
Sur la droite graduée, on le représente par un trait épais partant de $-3$ (exclu, cercle vide) vers la droite.
\end{exemplebox}

\begin{popfigure}
\begin{tikzpicture}[scale=1.5, >=Stealth]
    \draw[thick, <->] (-5, 0) -- (2, 0);
    \foreach \x in {-4,-3,-2,-1,0,1}
        \draw (\x, -0.1) -- (\x, 0.1) node[below=3pt] {\small $\x$};
    
    % Intervalle ]-3 ; +inf[
    \draw[popblue, very thick] (-3, 0) -- (1.5, 0);
    \draw[popblue, fill=white, thick] (-3, 0) circle (3pt); % cercle vide
    \node[popblue, above] at (-3, 0.2) {$-3$ (exclu)};
    \node[popblue, above] at (0, 0.4) {Ensemble des solutions $S = ]-3 ; +\infty[$};
\end{tikzpicture}
\legende{Représentation graphique de l'ensemble des solutions de $x > -3$.}
\end{popfigure}

\courssub{Méthode : Résoudre une inéquation du premier degré $ax + b < 0$ (ou $>$, $\le$, $\ge$)}

\begin{methode}[Résolution algorithmique]
Pour résoudre $ax + b < 0$ (avec $a \ne 0$) :
\begin{enumerate}
    \item \textbf{Isoler le terme en $x$} en ajoutant/retranchant $b$ des deux membres (Règle 1).
    \item \textbf{Diviser par $a$} (Règle 2 ou 3 selon le signe de $a$).
    \begin{itemize}
        \item Si $a > 0$ : le sens est conservé.
        \item Si $a < 0$ : \textbf{le sens est inversé}.
    \end{itemize}
    \item \textbf{Écrire l'ensemble des solutions $S$} sous forme d'intervalle ou d'inégalité.
    \item \textbf{Vérifier} (optionnel mais recommandé) en testant une valeur prise dans $S$ et une valeur prise hors de $S$.
\end{enumerate}
\end{methode}

\begin{exoresolu}[Résolution complète]
Résoudre l'inéquation : $\mathbf{5 - 2x \le 1}$.

\textbf{Solution détaillée :}
\begin{align*}
5 - 2x &\le 1 \\
-2x &\le 1 - 5 \quad \text{(on retranche 5 des deux membres, Règle 1)} \\
-2x &\le -4 \\
x &\ge \frac{-4}{-2} \quad \text{(on divise par -2, \textbf{négatif} : on \textbf{inverse} le sens, Règle 3)} \\
x &\ge 2
\end{align*}
L'ensemble des solutions est $S = [2 ; +\infty[$.

\textbf{Vérification :}
\begin{itemize}
    \item Test $x = 3 \in S$ : $5 - 2(3) = -1 \le 1$ \quad $\checkmark$ Vrai.
    \item Test $x = 0 \notin S$ : $5 - 0 = 5 \le 1$ \quad $\checkmark$ Faux.
\end{itemize}
\end{exoresolu}

\begin{savaistu}[Histoire des symboles : Thomas Harriot (1631)]
Les symboles $<$ et $>$ que nous utilisons tous les jours sont relativement récents ! Ils ont été introduits par le mathématicien et astronome anglais \textbf{Thomas Harriot} (1560--1621) dans son ouvrage \emph{Artis Analyticae Praxis}, publié à titre posthume en \textbf{1631}.
Avant lui, les mathématiciens écrivaient les inégalités en toutes lettres (ex: « $a$ est plus grand que $b$ ») ou utilisaient des symboles compliqués. Harriot a eu l'intuition géniale que la forme de la bouche grande ouverte ($>$) « mange » le plus grand nombre, tandis que la pointe ($<$) pointe vers le plus petit. Une mnémotechnique visuelle qui a traversé quatre siècles et qui est aujourd'hui universelle, de Yaoundé à New York en passant par Paris !
\end{savaistu}

\begin{exercice}[Entraînement immédiat]
Résoudre les inéquations suivantes et donner l'ensemble des solutions $S$.
\begin{enumerate}
    \item $3x - 7 > 2$
    \item $-4x + 1 \le 9$
    \item $5 - \frac{x}{2} < 1$
    \item $2(x - 3) \ge 4x + 1$
\end{enumerate}
\end{exercice}

\begin{corrige}[Exercice n°1]
\begin{enumerate}
    \item $3x - 7 > 2 \iff 3x > 9 \iff x > 3$. \quad $S = ]3 ; +\infty[$.
    \item $-4x + 1 \le 9 \iff -4x \le 8 \iff x \ge -2$ (\textbf{sens inversé} car division par $-4$). \quad $S = [-2 ; +\infty[$.
    \item $5 - \frac{x}{2} < 1 \iff -\frac{x}{2} < -4 \iff x > 8$ (\textbf{sens inversé} car multiplication par $-2$). \quad $S = ]8 ; +\infty[$.
    \item $2(x - 3) \ge 4x + 1 \iff 2x - 6 \ge 4x + 1 \iff -2x \ge 7 \iff x \le -\frac{7}{2}$ (\textbf{sens inversé}). \quad $S = ]-\infty ; \num{-3,5}]$.
\end{enumerate}
\end{corrige}

\coursec{Inéquations de la forme $ax + b > 0$}

Dans la section précédente, nous avons maîtrisé les \textbf{règles de transformation des inégalités} : on peut ajouter ou soustraire le même nombre aux deux membres sans changer le sens de l'inégalité, et on peut multiplier ou diviser les deux membres par un nombre \textbf{strictement positif} sans changer le sens ; en revanche, \textbf{multiplier ou diviser par un nombre strictement négatif inverse le sens de l'inégalité}. C'est exactement cet outil qui va nous permettre de résoudre toutes les inéquations du premier degré. Rappelons-nous la différence fondamentale : une équation demande « \emph{quelle valeur rend les deux membres égaux ?} » ; une inéquation demande « \emph{quelles valeurs rendent un membre plus grand (ou plus petit) que l'autre ?} ». La réponse ne sera donc plus un nombre unique, mais une \emph{infinité} de valeurs, que nous saurons décrire rigoureusement et représenter sur une droite graduée.

\courssub{Définition et intuition}

\begin{definition}[Inéquation du premier degré à une inconnue]
Une \cle{inéquation du premier degré à une inconnue} est une inéquation qui, après simplification des deux membres, peut s'écrire sous l'une des quatre formes suivantes, où $a$ et $b$ sont des nombres réels avec $a \neq 0$ :
\[
ax + b > 0 \quad ; \quad ax + b < 0 \quad ; \quad ax + b \geq 0 \quad ; \quad ax + b \leq 0.
\]
L'inconnue est généralement notée $x$. \cle{Résoudre} une inéquation, c'est trouver \textbf{toutes} les valeurs de $x$ qui rendent l'inégalité vraie. L'ensemble de ces valeurs est appelé \cle{ensemble des solutions} et est noté $S$.
\end{definition}

\begin{exemplebox}[Reconnaître une inéquation du premier degré]
Les inéquations suivantes sont du premier degré :
\begin{itemize}
    \item $3x - 5 > 0$ (déjà sous la forme $ax+b>0$, avec $a=3$ et $b=-5$).
    \item $2x + 1 < 4x - 7$ (on ramène tout dans le membre de gauche : $2x + 1 - 4x + 7 < 0$, soit $-2x + 8 < 0$).
    \item $\dfrac{x}{2} + 3 \geq 5$ (on multiplie les deux membres par $2 > 0$ : $x + 6 \geq 10$, puis $x - 4 \geq 0$).
\end{itemize}
En revanche, l'inéquation $x^2 - 4 > 0$ n'est \textbf{pas} du premier degré, car elle contient un terme en $x^2$ : elle ne peut pas se ramener à la forme $ax + b > 0$.
\end{exemplebox}

\emph{Intuition.} Résoudre $ax + b > 0$, c'est chercher tous les nombres $x$ pour lesquels l'expression $ax + b$ est \textbf{strictement positive}. Prenons l'exemple concret d'une vendeuse : si son bénéfice journalier est $50x - 5000$ (en FCFA) pour $x$ beignets vendus, se demander quand elle gagne de l'argent, c'est résoudre $50x - 5000 > 0$. La réponse ne sera pas « un nombre de beignets », mais « tous les nombres de beignets à partir d'un certain seuil » : une \textbf{demi-droite} de valeurs.

\courssub{Résolution de $ax + b > 0$ : discussion selon le signe de $a$}

C'est le cœur de la leçon. Nous allons transformer l'inéquation pour isoler $x$, exactement comme pour une équation, mais en \textbf{surveillant le signe de $a$} au moment de la division.

\begin{propriete}[Résolution de $ax + b > 0$]
Soient $a$ et $b$ deux nombres réels avec $a \neq 0$.
\begin{enumerate}
    \item \textbf{Si $a > 0$ :} on divise les deux membres par $a$ (positif), le sens de l'inégalité est \textbf{conservé}.
    \[
    ax + b > 0 \iff ax > -b \iff x > -\frac{b}{a}.
    \]
    L'ensemble des solutions est $S = \left] -\dfrac{b}{a} \ ; \ +\infty \right[$.
    \item \textbf{Si $a < 0$ :} on divise les deux membres par $a$ (négatif), le sens de l'inégalité est \textbf{inversé}.
    \[
    ax + b > 0 \iff ax > -b \iff x < -\frac{b}{a}.
    \]
    L'ensemble des solutions est $S = \left] -\infty \ ; \ -\dfrac{b}{a} \right[$.
\end{enumerate}
\end{propriete}

\begin{experience}[Démonstration guidée : pourquoi le signe de $a$ change tout ?]
Résolvons $-2x + 6 > 0$ de deux façons pour bien comprendre la règle.
\begin{enumerate}
    \item \textbf{Méthode 1 (division par $-2$) :}
    $-2x + 6 > 0 \iff -2x > -6$ (on soustrait $6$ aux deux membres, sens inchangé).
    On divise par $-2$, qui est \textbf{négatif} : \textbf{on inverse le sens} de l'inégalité : $x < \dfrac{-6}{-2}$, c'est-à-dire $x < 3$.
    \item \textbf{Méthode 2 (on évite de diviser par un négatif) :}
    $-2x + 6 > 0 \iff 6 > 2x$ (on a ajouté $2x$ aux deux membres, sens inchangé).
    $\iff 2x < 6$ (on écrit la même inégalité en échangeant les membres, ce qui inverse le sens).
    $\iff x < 3$ (on divise par $2 > 0$, sens inchangé).
\end{enumerate}
Les deux méthodes donnent le même résultat : $x < 3$. \textbf{Conclusion :} diviser les deux membres d'une inégalité par un nombre strictement négatif impose d'\textbf{inverser le symbole} d'inégalité ($>$ devient $<$, $\geq$ devient $\leq$, et réciproquement). C'est la seule différence avec la résolution d'une équation !
\end{experience}

\begin{propriete}[Tableau récapitulatif des quatre formes]
La méthode est identique pour les quatre formes : on conserve le type d'inégalité (stricte ou large) et on applique la règle du signe de $a$.
\begin{center}
\begin{tabular}{|c|c|c|}\hline
\rowcolor{popblueL}
\textbf{Inéquation} & \textbf{Cas $a > 0$} & \textbf{Cas $a < 0$} \\ \hline
$ax + b > 0$ & $x > -\dfrac{b}{a}$ & $x < -\dfrac{b}{a}$ \\ \hline
$ax + b < 0$ & $x < -\dfrac{b}{a}$ & $x > -\dfrac{b}{a}$ \\ \hline
$ax + b \geq 0$ & $x \geq -\dfrac{b}{a}$ & $x \leq -\dfrac{b}{a}$ \\ \hline
$ax + b \leq 0$ & $x \leq -\dfrac{b}{a}$ & $x \geq -\dfrac{b}{a}$ \\ \hline
\end{tabular}
\end{center}
\textbf{Remarque :} il est inutile d'apprendre ce tableau par cœur. Il suffit de retenir une seule règle : \emph{quand on multiplie ou on divise par un nombre strictement négatif, on inverse le sens de l'inégalité}.
\end{propriete}

\courssub{Représentation graphique et écriture de l'ensemble des solutions}

Une fois l'ensemble solution $S$ trouvé (par exemple $x > 3$ ou $x \leq -2$), il faut savoir le \textbf{représenter sur une droite graduée} : c'est une compétence exigible au BEPC.

\begin{methode}[Représenter l'ensemble des solutions sur une droite graduée]
\begin{enumerate}
    \item Tracer une droite horizontale graduée et placer la valeur limite $-\dfrac{b}{a}$.
    \item \textbf{Crochet (ou cercle) OUVERT :} si l'inégalité est \textbf{stricte} ($>$ ou $<$). La valeur limite n'est \textbf{pas} solution : on dessine un petit cercle vide.
    \item \textbf{Crochet (ou point) FERMÉ :} si l'inégalité est \textbf{large} ($\geq$ ou $\leq$). La valeur limite \textbf{est} solution : on dessine un point plein.
    \item \textbf{Colorier (ou hachurer) la partie solution :}
    \begin{itemize}
        \item $x > k$ ou $x \geq k$ : on colorie vers la \textbf{droite} (vers les grandes valeurs).
        \item $x < k$ ou $x \leq k$ : on colorie vers la \textbf{gauche} (vers les petites valeurs).
    \end{itemize}
\end{enumerate}
\end{methode}

\begin{popfigure}
\begin{tikzpicture}[scale=0.9, every node/.style={font=\small}]
    % Axe principal
    \draw[->, thick, popdark] (-1,0) -- (7,0) node[right] {$x$};
    % Graduations
    \foreach \x in {0,1,2,3,4,5,6} {
        \draw[thick] (\x, -0.1) -- (\x, 0.1) node[below=2pt] {$\x$};
    }
    % Solution x > 3 : cercle vide en 3, coloriage vers la droite
    \draw[popblue, very thick] (3, 0) -- (6.6, 0);
    \draw[popblue, very thick, ->] (6.6, 0) -- (6.9, 0);
    \draw[popblue, thick, fill=white] (3,0) circle (3pt);
    \node[popblue, above] at (3, 0.4) {$x > 3$};
    \node[popblue, above] at (5, 0.15) {$S = \left] 3 \ ; \ +\infty \right[$};
\end{tikzpicture}
\legende{Représentation des solutions de $2x - 6 > 0$, c'est-à-dire $x > 3$ : cercle vide en $3$ (la valeur $3$ n'est pas solution), coloriage vers la droite.}
\end{popfigure}

\begin{popfigure}
\begin{tikzpicture}[scale=0.9, every node/.style={font=\small}]
    % Axe principal
    \draw[->, thick, popdark] (-1,0) -- (7,0) node[right] {$x$};
    % Graduations
    \foreach \x in {0,1,2,3,4,5,6} {
        \draw[thick] (\x, -0.1) -- (\x, 0.1) node[below=2pt] {$\x$};
    }
    % Solution x <= 3 : point plein en 3, coloriage vers la gauche
    \draw[poporange, very thick] (0.4, 0) -- (3, 0);
    \draw[poporange, very thick, <-] (0.1, 0) -- (0.4, 0);
    \fill[poporange] (3,0) circle (3pt);
    \node[poporange, above] at (3, 0.4) {$x \leq 3$};
    \node[poporange, above] at (1.5, 0.15) {$S = \left] -\infty \ ; \ 3 \right]$};
\end{tikzpicture}
\legende{Représentation des solutions de $-3x + 9 \geq 0$, c'est-à-dire $x \leq 3$ : point plein en $3$ (la valeur $3$ est solution), coloriage vers la gauche.}
\end{popfigure}

\begin{aretenir}[Écriture de l'ensemble des solutions]
En classe de 3ème, l'ensemble des solutions peut s'écrire de deux manières :
\begin{itemize}
    \item \textbf{En langage ensembliste :} par exemple $S = \{ x \in \mathbb{R} \mid x > 3 \}$, qui se lit « ensemble des réels $x$ tels que $x > 3$ ».
    \item \textbf{En notation par intervalles} (très utilisée ensuite au lycée) :
    \begin{itemize}
        \item $x > a \iff S = \left] a \ ; \ +\infty \right[$ \quad (crochet ouvert en $a$)
        \item $x \geq a \iff S = \left[ a \ ; \ +\infty \right[$ \quad (crochet fermé en $a$)
        \item $x < a \iff S = \left] -\infty \ ; \ a \right[$ \quad (crochet ouvert en $a$)
        \item $x \leq a \iff S = \left] -\infty \ ; \ a \right]$ \quad (crochet fermé en $a$)
    \end{itemize}
\end{itemize}
\textbf{Règle mnémotechnique :} le crochet est \textbf{toujours ouvert} du côté de l'infini ($-\infty$ ou $+\infty$), car l'infini n'est pas un nombre que l'on peut atteindre. Du côté de la valeur $a$, le crochet suit l'inégalité : \textbf{ouvert} pour une inégalité stricte ($>$ ou $<$), \textbf{fermé} pour une inégalité large ($\geq$ ou $\leq$).
\end{aretenir}

\courssub{Vérification des solutions}

Pour une équation, on vérifie en remplaçant $x$ par la solution unique. Pour une inéquation, il y a une infinité de solutions : on ne peut pas toutes les tester ! On utilise alors la \textbf{méthode des valeurs test}.

\begin{methode}[Vérifier la résolution d'une inéquation]
\begin{enumerate}
    \item Choisir une valeur \textbf{à l'intérieur} de l'ensemble solution trouvé (une valeur simple, par exemple $0$ si elle convient).
    \item Choisir une valeur \textbf{à l'extérieur} de l'ensemble solution (de l'autre côté de la borne).
    \item Remplacer $x$ par chacune de ces deux valeurs dans l'inéquation \textbf{d'origine}.
    \item La première valeur doit rendre l'inéquation \textbf{vraie}, la seconde doit la rendre \textbf{fausse}. Si c'est le cas, la résolution est très probablement correcte.
\end{enumerate}
\end{methode}

\begin{exemplebox}[Vérification pour $-2x + 8 > 0$]
Nous avons trouvé $S = \left] -\infty \ ; \ 4 \right[$, c'est-à-dire $x < 4$.
\begin{itemize}
    \item \textbf{Valeur solution :} prenons $x = 0$ (on a bien $0 < 4$).
    $-2 \times 0 + 8 = 8$, et $8 > 0$ : l'inégalité est \textbf{vraie}. $\checkmark$
    \item \textbf{Valeur non solution :} prenons $x = 5$ (on a $5 > 4$).
    $-2 \times 5 + 8 = -10 + 8 = -2$, et $-2 > 0$ est \textbf{faux}. $\checkmark$
\end{itemize}
Les deux tests sont conformes : la résolution est validée.
\end{exemplebox}

\courssub{Exemples résolus et pièges à éviter}

\begin{exoresolu}[Résolution complète avec $a < 0$]
\textbf{Énoncé :} Résoudre dans $\mathbb{R}$ l'inéquation $-2x + 8 > 0$, puis représenter l'ensemble des solutions sur une droite graduée.
\tcblower
\textbf{Solution détaillée :}
\begin{align*}
-2x + 8 &> 0 \\
-2x &> -8 \quad \text{(on soustrait } 8 \text{ aux deux membres, sens inchangé)} \\
x &< \frac{-8}{-2} \quad \text{(on divise par } -2 < 0 \text{ : \textbf{on inverse le sens})} \\
x &< 4
\end{align*}
\textbf{Ensemble des solutions :} $S = \{ x \in \mathbb{R} \mid x < 4 \} = \left] -\infty \ ; \ 4 \right[$.

\textbf{Représentation graphique :}
\begin{center}
\begin{tikzpicture}[scale=0.8]
    \draw[->, thick, popdark] (-1,0) -- (6,0) node[right] {$x$};
    \foreach \x in {0,1,2,3,4,5} \draw (\x, -0.1) -- (\x, 0.1) node[below] {$\x$};
    \draw[popblue, very thick] (0.4, 0) -- (4, 0);
    \draw[popblue, very thick, <-] (0.1, 0) -- (0.4, 0);
    \draw[popblue, thick, fill=white] (4,0) circle (3pt);
    \node[popblue, above] at (4, 0.4) {$x < 4$};
\end{tikzpicture}
\end{center}
Cercle vide en $4$ (inégalité stricte), coloriage vers la gauche.
\end{exoresolu}

\begin{exoresolu}[Inéquation à coefficients fractionnaires]
\textbf{Énoncé :} Résoudre dans $\mathbb{R}$ l'inéquation $\dfrac{x}{3} - 2 \leq 1$.
\tcblower
\textbf{Solution détaillée :}
\begin{align*}
\frac{x}{3} - 2 &\leq 1 \\
\frac{x}{3} &\leq 3 \quad \text{(on ajoute } 2 \text{ aux deux membres, sens inchangé)} \\
x &\leq 9 \quad \text{(on multiplie par } 3 > 0 \text{, sens inchangé)}
\end{align*}
\textbf{Ensemble des solutions :} $S = \left] -\infty \ ; \ 9 \right]$. Point \textbf{plein} en $9$ (inégalité large), coloriage vers la gauche.
\end{exoresolu}

\begin{exoresolu}[Inéquation avec des termes en $x$ dans les deux membres]
\textbf{Énoncé :} Résoudre dans $\mathbb{R}$ l'inéquation $5x - 3 \geq 7x + 5$.
\tcblower
\textbf{Solution détaillée :} on regroupe d'abord les termes en $x$ dans un membre et les constantes dans l'autre.
\begin{align*}
5x - 3 &\geq 7x + 5 \\
5x - 7x &\geq 5 + 3 \quad \text{(on soustrait } 7x \text{ et on ajoute } 3 \text{ aux deux membres)} \\
-2x &\geq 8 \\
x &\leq \frac{8}{-2} \quad \text{(division par } -2 < 0 \text{ : \textbf{on inverse le sens})} \\
x &\leq -4
\end{align*}
\textbf{Ensemble des solutions :} $S = \left] -\infty \ ; \ -4 \right]$.
\textbf{Vérification :} pour $x = -5$ (solution) : $5 \times (-5) - 3 = -28$ et $7 \times (-5) + 5 = -30$ ; or $-28 \geq -30$ est vrai. $\checkmark$ Pour $x = 0$ (non solution) : $-3 \geq 5$ est faux. $\checkmark$
\end{exoresolu}

\begin{attention}[Erreurs fréquentes — Pièges du BEPC]
\begin{enumerate}
    \item \textbf{Oublier d'inverser le sens en divisant par un nombre négatif.} C'est l'erreur la plus fréquente.
    \emph{Exemple :} de $-3x > 6$, conclure $x > -2$ est \textbf{faux}. La bonne conclusion est $x < -2$.
    \item \textbf{Confondre cercle vide et point plein sur la droite graduée.}
    $x > 3$ : cercle \textbf{vide} en $3$. \quad $x \geq 3$ : point \textbf{plein} en $3$.
    \item \textbf{Donner une seule valeur comme solution.} Écrire « la solution est $x = 4$ » pour une inéquation est une erreur grave : la solution est un \textbf{ensemble de valeurs} (une demi-droite), pas un nombre.
    \item \textbf{Mal écrire l'intervalle :} écrire $\left[ 3 \ ; \ +\infty \right[$ pour $x > 3$ (crochet fermé au lieu d'ouvert). Relire la règle : « strict $\to$ ouvert, large $\to$ fermé ».
    \item \textbf{Ne pas aller au bout de la simplification.} Certains cas particuliers peuvent survenir : $2x + 4 > 2x + 1 \iff 4 > 1$, ce qui est \textbf{toujours vrai}, donc $S = \mathbb{R}$. De même, $2x + 4 < 2x + 1 \iff 4 < 1$, ce qui est \textbf{toujours faux}, donc $S = \varnothing$ (aucune solution).
\end{enumerate}
\end{attention}

\courssub{Applications concrètes : seuils et marges}

\begin{savaistu}[Les inéquations au quotidien : seuils de rentabilité et marges de sécurité]
Les inéquations ne sont pas des exercices abstraits : ce sont les outils mathématiques de la \textbf{prise de décision} quand on cherche une limite à ne pas dépasser ou un seuil à atteindre.
\begin{itemize}
    \item \textbf{Commerçante (marché Central de Yaoundé) :} une vendeuse de beignets a un coût fixe de 5\,000 FCFA par jour (gaz, transport) et un coût variable de 50 FCFA par beignet. Elle vend chaque beignet 100 FCFA. Son bénéfice pour $x$ beignets vendus est $B(x) = 100x - (50x + 5000) = 50x - 5000$. Pour être \textbf{bénéficiaire}, il faut $B(x) > 0$, c'est-à-dire $50x - 5000 > 0 \iff x > 100$. Elle doit donc vendre \textbf{au moins 101 beignets} par jour pour gagner de l'argent : c'est son \textbf{seuil de rentabilité}.
    \item \textbf{Ingénieur (ouvrages comme le barrage de Lom Pangar) :} un câble supporte une tension maximale de 10\,000 N. La tension réelle est $T = 200x + 500$ (en newtons), où $x$ est la charge suspendue (en kg). Pour la sécurité, on impose $T \leq 10\,000$, soit $200x + 500 \leq 10\,000 \iff 200x \leq 9\,500 \iff x \leq 47,5$. La charge maximale autorisée est donc $47,5$ kg : c'est une \textbf{marge de sécurité}.
    \item \textbf{Électrification rurale (ENEO) :} pour dimensionner un transformateur, on exige que la puissance demandée $P$ reste \textbf{inférieure} à $80\,\%$ de la puissance nominale $P_n$ : $P \leq 0,8\,P_n$. Cette inéquation guide le choix du matériel à installer dans un village.
\end{itemize}
Dans tous ces cas, la réponse n'est pas un nombre unique, mais une \textbf{plage de valeurs acceptables}. C'est toute la puissance des inéquations.
\end{savaistu}

\begin{exercice}[Entraînement — À toi de jouer !]
\begin{enumerate}
    \item Résoudre dans $\mathbb{R}$ et représenter les solutions sur une droite graduée :
    \begin{enumerate}
        \item $5x - 15 \geq 0$
        \item $-4x - 8 < 0$
        \item $3 - 2x \leq 7$
        \item $\dfrac{x}{5} + 1 > 3$
    \end{enumerate}
    \item Un taxi-man à Yaoundé paie 3\,000 FCFA de location journalière pour son véhicule. Il gagne en moyenne 500 FCFA par course. Combien de courses, au minimum, doit-il faire dans la journée pour réaliser un bénéfice strictement positif ? Modéliser par une inéquation, puis résoudre.
    \item Vérifier, par la méthode des valeurs test, que l'ensemble des solutions de $-3x + 12 > 0$ est bien donné par $x < 4$.
\end{enumerate}
\end{exercice}

\begin{corrige}[Exercice n°1]
\begin{enumerate}
    \item \textbf{a)} $5x - 15 \geq 0 \iff 5x \geq 15 \iff x \geq 3$ (division par $5 > 0$, sens inchangé). $S = \left[ 3 \ ; \ +\infty \right[$. Point \textbf{plein} en $3$, coloriage vers la droite.

    \textbf{b)} $-4x - 8 < 0 \iff -4x < 8 \iff x > -2$ (division par $-4 < 0$, \textbf{inversion} du sens). $S = \left] -2 \ ; \ +\infty \right[$. Cercle \textbf{vide} en $-2$, coloriage vers la droite.

    \textbf{c)} $3 - 2x \leq 7 \iff -2x \leq 4 \iff x \geq -2$ (division par $-2 < 0$, \textbf{inversion} du sens). $S = \left[ -2 \ ; \ +\infty \right[$. Point \textbf{plein} en $-2$, coloriage vers la droite.

    \textbf{d)} $\dfrac{x}{5} + 1 > 3 \iff \dfrac{x}{5} > 2 \iff x > 10$ (multiplication par $5 > 0$, sens inchangé). $S = \left] 10 \ ; \ +\infty \right[$. Cercle \textbf{vide} en $10$, coloriage vers la droite.

    \item Soit $x$ le nombre de courses. Le bénéfice journalier est $B(x) = 500x - 3000$ (en FCFA). On veut $B(x) > 0$ :
    \[
    500x - 3000 > 0 \iff 500x > 3000 \iff x > 6.
    \]
    Comme $x$ est un nombre entier de courses et que $x > 6$, le taxi-man doit faire \textbf{au moins 7 courses} dans la journée.

    \item Pour $-3x + 12 > 0$, la résolution donne $-3x > -12 \iff x < 4$ (inversion du sens, division par $-3 < 0$).
    \textbf{Vérification :}
    \begin{itemize}
        \item Valeur solution : $x = 0$ (on a $0 < 4$) : $-3 \times 0 + 12 = 12$, et $12 > 0$ est \textbf{vrai}. $\checkmark$
        \item Valeur non solution : $x = 5$ (on a $5 > 4$) : $-3 \times 5 + 12 = -15 + 12 = -3$, et $-3 > 0$ est \textbf{faux}. $\checkmark$
    \end{itemize}
    Les deux tests sont conformes : la résolution est correcte.
\end{enumerate}
\end{corrige}

\coursec{Résolution de problèmes concrets}

Après avoir maîtrisé la résolution des équations et inéquations du premier degré à une inconnue dans $\mathbb{R}$, nous allons voir comment ces outils permettent de modéliser et de résoudre des situations concrètes : comparaison de forfaits téléphoniques, calcul de seuils de rentabilité, problèmes géométriques soumis à des contraintes, etc. La démarche reste toujours la même : \textbf{traduire} le problème en langage mathématique, \textbf{résoudre}, puis \textbf{interpréter} la solution dans le contexte initial.

\courssub{Méthodologie générale : les 5 étapes}

\begin{methode}[Résolution d'un problème concret par équation ou inéquation]
\begin{enumerate}
\item \textbf{Choisir l'inconnue} : identifier la quantité recherchée et la noter par une lettre (souvent $x$). Préciser son unité et son \cle{ensemble de définition} (entier naturel $\mathbb{N}$, réel positif $\mathbb{R}_+^*$, etc.).
\item \textbf{Traduire en équation ou inéquation} : exprimer les relations du problème à l'aide des opérations arithmétiques. On obtient une \cle{équation} ($=$) si on cherche une valeur exacte (seuil), une \cle{inéquation} ($<$, $>$, $\le$, $\ge$) si on cherche un intervalle de valeurs.
\item \textbf{Résoudre} : appliquer les règles de transformation vues précédemment (addition, soustraction, multiplication, division en respectant le signe du coefficient). Pour une inéquation, ne pas oublier d'\textbf{inverser le sens} si on divise ou multiplie par un nombre \textbf{négatif}.
\item \textbf{Interpréter la solution mathématique} : revenir au contexte. La solution est-elle un nombre entier ? Doit-elle être positive ? Est-elle bornée par des contraintes physiques (ex. nombre de minutes, longueur d'un côté) ?
\item \textbf{Conclure par une phrase complète} : rédiger la réponse en français, en justifiant chaque étape et en vérifiant la cohérence avec l'énoncé.
\end{enumerate}
\end{methode}

\courssub{Problèmes de seuil et comparaison d'offres}

Les situations de \cle{seuil} apparaissent dès qu'on compare deux formules de coût : l'une contient une partie fixe et une partie variable, l'autre est purement variable. Le seuil correspond à la valeur de l'inconnue pour laquelle les deux coûts sont égaux ; au-delà, l'une des offres devient plus avantageuse.

\begin{exemplebox}[Comparaison de deux forfaits téléphoniques]
Soit deux forfaits proposés par MTN et Orange :
\begin{itemize}
\item Forfait A : 1\,000 FCFA d'abonnement mensuel + 10 FCFA par minute de communication.
\item Forfait B : 25 FCFA par minute, sans abonnement.
\end{itemize}
On cherche à partir de quel nombre de minutes le forfait A devient moins cher que le forfait B.
\end{exemplebox}

\noindent\textbf{Modélisation.} Soit $x$ le nombre de minutes consommées dans le mois ($x\in\mathbb{N}$). Le coût du forfait A est $C_A(x)=1000+10x$ (en FCFA), celui du forfait B est $C_B(x)=25x$ (en FCFA). On cherche $x$ tel que $C_A(x) < C_B(x)$, c'est-à-dire
\[
1000+10x < 25x \quad\Longleftrightarrow\quad 1000 < 15x \quad\Longleftrightarrow\quad x > \frac{1000}{15}=66,\overline{6}.
\]
Comme $x$ est un nombre entier de minutes, le plus petit entier strictement supérieur à $66,\overline{6}$ est $67$. Donc \textbf{à partir de 67 minutes, le forfait A est plus avantageux}.

\begin{popfigure}
\begin{tikzpicture}
\begin{axis}[
    width=\linewidth,
    height=0.55\linewidth,
    axis lines=middle,
    xlabel={Minutes $x$},
    ylabel={Coût (FCFA)},
    xmin=0, xmax=120,
    ymin=0, ymax=3500,
    xtick={0,20,40,60,66.66,80,100,120},
    xticklabels={0,20,40,60,{66,7},80,100,120},
    ytick={0,500,1000,1500,2000,2500,3000,3500},
    legend pos=north west,
    grid=major,
    grid style={popline},
    tick label style={font=\small},
    label style={font=\small},
]
\addplot[popblue, thick, domain=0:120, samples=2] {1000+10*x};
\addlegendentry{Forfait A : $1000+10x$}
\addplot[poporange, thick, domain=0:120, samples=2] {25*x};
\addlegendentry{Forfait B : $25x$}
% Intersection point
\addplot[only marks, mark=*, mark size=2.5, popgreen] coordinates {(66.666,2666.66)};
\node[above right, font=\small, popgreen] at (axis cs:66.666,2666.66) {Seuil $\approx 66,7$ min};
\end{axis}
\end{tikzpicture}
\legende{Comparaison graphique des deux forfaits : le point d'intersection donne le seuil à partir duquel le forfait A devient moins cher.}
\end{popfigure}

\begin{attention}[Piège fréquent : réponse mathématique vs réponse contextuelle]
Dans l'exemple ci-dessus, la résolution de l'inéquation donne $x > 66,\overline{6}$. Si l'on répond « $x > 66,\overline{6}$ minutes » sans tenir compte du fait que le nombre de minutes est un \cle{entier naturel}, on commet une erreur de rédaction. De même, pour un problème de fabrication d'objets (beignets, briques, etc.), une solution $x=7,5$ est mathématiquement correcte pour l'équation mais \textbf{impossible dans la réalité} : il faut arrondir à l'entier supérieur ou inférieur selon le sens de l'inégalité et expliciter ce choix dans la conclusion.
\end{attention}

\courssub{Problèmes géométriques simples}

Les contraintes sur un périmètre, une aire ou une longueur se traduisent souvent par une inéquation. La variable représente alors une dimension (longueur, largeur, rayon) qui doit rester \textbf{strictement positive} (appartenance à $\mathbb{R}_+^*$).

\begin{exemplebox}[Rectangle de périmètre limité]
Un agriculteur veut clôturer un rectangle dont la longueur est supérieure de 5 m à la largeur. Il dispose de 80 m de grillage. Quelles sont les dimensions possibles ?
\end{exemplebox}

\noindent Soit $x$ la largeur en mètres ($x>0$, soit $x\in\mathbb{R}_+^*$). La longueur est $x+5$. Le périmètre est $P=2(x+x+5)=4x+10$. La contrainte « au plus 80 m » donne l'inéquation
\[
4x+10 \le 80 \;\Longrightarrow\; 4x \le 70 \;\Longrightarrow\; x \le 17,5.
\]
En tenant compte de la condition de positivité $x>0$, les largeurs possibles sont $0 < x \le 17,5$ m. La longueur correspondante varie entre $5$ m et $22,5$ m. On notera que $x$ peut être un nombre décimal (mesure continue), contrairement au nombre de minutes qui est discret.

\courssub{Validation de la réponse et rédaction complète}

Une fois la solution numérique obtenue, il est indispensable de :
\begin{itemize}
\item vérifier qu'elle appartient bien à l'ensemble de définition de l'inconnue (entier, positif, borné) ;
\item la tester dans l'énoncé original (remplacer $x$ par la valeur trouvée) ;
\item rédiger une conclusion sous forme de phrase affirmative, en reprenant les unités et le vocabulaire du problème.
\end{itemize}

\begin{exoresolu}[Problème de Mama Ngo Bell]
\textbf{Énoncé.} Mama Ngo Bell vend des beignets à 25 FCFA l'unité. Elle a dépensé 1\,500 FCFA pour les ingrédients. Combien de beignets doit-elle vendre pour (a) couvrir exactement ses frais ? (b) réaliser un bénéfice ?
\textbf{Solution.}
\begin{enumerate}
\item Soit $x$ le nombre de beignets vendus ($x\in\mathbb{N}$). Le chiffre d'affaires est $25x$ FCFA.
\item (a) Couvrir les frais (seuil de rentabilité) : $25x = 1500 \;\Rightarrow\; x = 60$. Elle doit vendre exactement 60 beignets.
\item (b) Faire un bénéfice : $25x > 1500 \;\Rightarrow\; x > 60$. Comme $x$ est entier, le plus petit entier strictement supérieur à 60 est 61. Elle doit vendre \textbf{au moins 61 beignets}.
\end{enumerate}
\textbf{Conclusion.} Pour ne pas perdre d'argent, Mama Ngo Bell doit vendre 60 beignets ; pour gagner de l'argent, elle doit en vendre 61 ou plus.
\tcblower
\textbf{Analyse.} On a d'abord résolu une équation (seuil de rentabilité), puis une inéquation (zone de bénéfice). La transition de l'égalité à l'inégalité stricte illustre le passage du « point mort » à la « zone de profit ».
\end{exoresolu}

\courssub{Le savais-tu ?}

\begin{savaistu}[Modélisation et recherche opérationnelle]
La modélisation par des inéquations (et plus généralement par des systèmes d'inéquations) est au cœur de la \emph{recherche opérationnelle}. Les entreprises de transport de Yaoundé et Douala (SOTUC, taxis-brousses, camions de livraison) l'utilisent pour optimiser leurs tournées : minimiser la consommation de carburant, respecter les créneaux horaires, maximiser le chargement. Les logiciels de gestion de flotte résolvent en temps réel des milliers d'inéquations pour proposer les itinéraires les plus rentables !
\end{savaistu}

\courssub{Exercices d'application}

\begin{exercice}[Forfaits internet]
Un fournisseur propose deux offres :
\begin{itemize}
\item Offre 1 : 5\,000 FCFA/mois + 2 FCFA/Mo de données.
\item Offre 2 : 15\,000 FCFA/mois, données illimitées.
\end{itemize}
À partir de quel volume de données l'offre 2 devient-elle plus avantageuse ? Résolvez par inéquation et illustrez par un graphique.
\end{exercice}

\begin{exercice}[Clôture d'un terrain]
Un terrain rectangulaire a une longueur double de sa largeur. Le propriétaire dispose de 120 m de clôture. Déterminez l'intervalle des largeurs possibles.
\end{exercice}

\begin{corrige}[Exercice 1]
Soit $x$ le volume en Mo. Coût offre 1 : $5000+2x$ FCFA. Coût offre 2 : $15000$ FCFA. L'offre 2 est plus avantageuse si $5000+2x > 15000 \Rightarrow 2x > 10000 \Rightarrow x > 5000$. Donc pour plus de 5\,000 Mo (soit 5 Go), l'offre 2 est moins chère.
\end{corrige}

\begin{corrige}[Exercice 2]
Soit $x$ la largeur en mètres ($x>0$). Longueur $=2x$. Périmètre $=2(x+2x)=6x$. Contrainte : $6x \le 120 \Rightarrow x \le 20$. Comme $x>0$, on a $0 < x \le 20$ m.
\end{corrige}

\newpage

\coursec{Activité d'intégration}

\courssub{Objectifs de l'activité}
Cette activité d'intégration vise à mobiliser l'ensemble des savoirs et savoir-faire de l'unité « Équations et inéquations du premier degré à une inconnue dans $\mathbb{R}$ » dans une situation concrète de la vie économique camerounaise. Elle permet à l'élève de :
\begin{itemize}
    \item Traduire un problème réel en langage algébrique (choix de l'inconnue, écriture d'une équation et d'inéquations).
    \item Résoudre rigoureusement ces équations et inéquations en justifiant chaque étape (règles de transformation).
    \item Interpréter les solutions dans le contexte du problème (seuil de rentabilité, bénéfice minimal, comparaison de deux offres).
    \item Rédiger une démarche complète, structurée et argumentée, puis la communiquer oralement.
    \item Développer un esprit critique sur la modélisation mathématique : limites du modèle, hypothèses implicites (vente de la totalité de la production, prix constants).
\end{itemize}

\courssub{Situation}
\textbf{Contexte :} Au quartier Nkol-Eton à Yaoundé, comme dans de nombreux carrefours animés du Cameroun, le commerce informel de la « soya » (brochettes de viande grillée) bat son plein chaque soir. Deux vendeurs, \textbf{Monsieur Atangana} et \textbf{Monsieur Bello}, s'installent côte à côte au même carrefour. Leur stratégie d'approvisionnement et de fixation des prix diffère.

\textbf{Données pour Monsieur Atangana :}
\begin{itemize}
    \item Investissement initial (capital de départ pour le charbon, la viande, les épices, le transport) : \textbf{8\,000 FCFA}.
    \item Prix de vente unitaire d'une brochette : \textbf{200 FCFA}.
\end{itemize}

\textbf{Données pour Monsieur Bello :}
\begin{itemize}
    \item Investissement initial : \textbf{5\,000 FCFA} (il achète des morceaux moins chers et utilise un charbon de qualité inférieure).
    \item Prix de vente unitaire d'une brochette : \textbf{150 FCFA}.
\end{itemize}

\textbf{Hypothèses de modélisation (simplifications acceptées pour l'exercice) :}
\begin{enumerate}
    \item Toute la production du soir est vendue (pas d'invendus).
    \item Les prix de vente et les coûts fixes restent constants pendant la soirée.
    \item On ne considère que les coûts fixes initiaux (pas de coûts variables par brochette comme l'huile ou le piment, pour rester dans le cadre du 1er degré).
    \item Le nombre de brochettes vendues est un nombre entier naturel, mais on raisonnera dans $\mathbb{R}$ pour la résolution algébrique, puis on interprétera dans $\mathbb{N}$.
\end{enumerate}

\courssub{Consignes}
Travailler par groupes de 3 ou 4 élèves. Chaque groupe désigne un rapporteur qui présentera la démarche au tableau. Rédiger les réponses sur une feuille A4 en suivant scrupuleusement les étapes ci-dessous.

\begin{enumerate}
    \item \textbf{Modélisation pour Monsieur Atangana (Seuil de rentabilité).}
    \begin{enumerate}
        \item Choisir une inconnue $x$ adaptée et expliciter son sens (unité comprise).
        \item Exprimer le chiffre d'affaires $CA_A(x)$ en fonction de $x$.
        \item Écrire l'équation qui traduit : « Le chiffre d'affaires couvre exactement l'investissement initial ».
        \item Résoudre cette équation dans $\mathbb{R}$ en justifiant chaque passage (règle utilisée).
        \item Interpréter le résultat dans $\mathbb{N}$ : combien de brochettes minimum doit-il vendre pour ne pas perdre d'argent ? On note ce nombre $N_A$.
    \end{enumerate}

    \item \textbf{Modélisation pour Monsieur Atangana (Objectif de bénéfice).}
    \begin{enumerate}
        \item Écrire l'inéquation qui traduit : « Le bénéfice est \emph{au moins} de 5\,000 FCFA ».
        \item Rappeler la définition du bénéfice : $B(x) = CA(x) - \text{Investissement}$.
        \item Résoudre l'inéquation dans $\mathbb{R}$ en justifiant les règles de transformation (attention au signe du coefficient de $x$).
        \item Donner l'ensemble solution $S_A$ dans $\mathbb{R}$ et l'interpréter dans $\mathbb{N}$ : nombre minimal de brochettes pour atteindre cet objectif.
    \end{enumerate}

    \item \textbf{Comparaison des deux vendeurs (Analyse comparative).}
    \begin{enumerate}
        \item Écrire la fonction bénéfice de Monsieur Bello : $B_B(x) = 150x - 5000$ (où $x$ est le nombre de brochettes vendues par Bello).
        \item On suppose que les deux vendeurs vendent \emph{exactement le même nombre} $x$ de brochettes (même affluence au carrefour).
        \item Écrire et résoudre l'inéquation : « Le bénéfice de Monsieur Atangana est \emph{strictement supérieur} à celui de Monsieur Bello ».
        \item Interpréter : à partir de quel nombre de brochettes vendues la stratégie « investissement élevé / prix élevé » devient-elle plus rentable que la stratégie « investissement faible / prix bas » ?
    \end{enumerate}

    \item \textbf{Synthèse et communication.}
    \begin{enumerate}
        \item Rédiger une conclusion unique résumant les trois seuils trouvés ($N_A$, seuil bénéfice 5000, seuil supériorité sur Bello).
        \item Préparer l'oral : le rapporteur doit expliquer le choix de l'inconnue, les règles de transformation utilisées pour l'inéquation (question 2.c) et l'interprétation économique du dernier résultat.
    \end{enumerate}
\end{enumerate}

\courssub{Ressources à mobiliser}
\begin{aretenir}[Rappels de cours — Unité : Équations et Inéquations 1er degré]
\textbf{1. Équation de la forme $ax + b = 0$ ($a \neq 0$) :}
La solution unique est $x = -\dfrac{b}{a}$.
\emph{Règles de transformation (équivalentes) :}
\begin{itemize}
    \item On peut ajouter ou soustraire le même nombre (ou expression) des deux membres.
    \item On peut multiplier ou diviser les deux membres par un même nombre \textbf{non nul}.
\end{itemize}

\textbf{2. Inéquation de la forme $ax + b > 0$ (ou $\geq, <, \leq$) :}
\emph{Règles de transformation :}
\begin{itemize}
    \item On peut ajouter ou soustraire le même nombre des deux membres \textbf{sans changer le sens de l'inégalité}.
    \item On peut multiplier ou diviser les deux membres par un même nombre \textbf{strictement positif} \textbf{sans changer le sens}.
    \item On peut multiplier ou diviser les deux membres par un même nombre \textbf{strictement négatif} \textbf{en inversant le sens de l'inégalité} (c'est le piège classique !).
\end{itemize}
La solution est un intervalle de $\mathbb{R}$ (ex: $]k; +\infty[$, $[k; +\infty[$, $]-\infty; k[$, $]-\infty; k]$).

\textbf{3. Vocabulaire économique simplifié :}
\begin{itemize}
    \item Chiffre d'affaires ($CA$) = Prix unitaire $\times$ Quantité vendue.
    \item Bénéfice ($B$) = Chiffre d'affaires $-$ Coûts (Investissement).
    \item Seuil de rentabilité (point mort) : $CA = \text{Investissement}$ $\Leftrightarrow B = 0$.
\end{itemize}
\end{aretenir}

\courssub{Démarche et corrigé commenté}
\begin{exoresolu}[Corrigé détaillé et commenté de l'activité d'intégration]
\textbf{Groupe : [Noms des élèves] \hfill Date : [Date du jour]}

\medskip
\textbf{\underline{1. Modélisation pour Monsieur Atangana — Seuil de rentabilité}}

\noindent \textbf{a) Choix de l'inconnue.}
\emph{Soit $x$ le nombre de brochettes vendues par Monsieur Atangana pendant la soirée.} $x \in \mathbb{N}$ (mais on résout dans $\mathbb{R}$).

\noindent \textbf{b) Chiffre d'affaires.}
Chaque brochette est vendue 200 FCFA.
\[ CA_A(x) = 200 \times x = 200x \quad \text{(en FCFA)} \]

\noindent \textbf{c) Équation du seuil de rentabilité.}
« Couvrir exactement l'investissement » signifie $CA_A(x) = 8000$.
\[ 200x = 8000 \]

\noindent \textbf{d) Résolution dans $\mathbb{R}$.}
\begin{align*}
200x &= 8000 \\
\dfrac{200x}{200} &= \dfrac{8000}{200} \qquad \text{(Division des deux membres par 200, nombre non nul, équivalence conservée)} \\
x &= 40
\end{align*}
L'ensemble solution dans $\mathbb{R}$ est $S = \{40\}$.

\noindent \textbf{e) Interprétation dans $\mathbb{N}$.}
$N_A = 40$. Monsieur Atangana doit vendre \textbf{au moins 40 brochettes} pour rembourser ses 8\,000 FCFA. S'il en vend 39, il perd de l'argent ; à 40, il est à l'équilibre (bénéfice nul).

\medskip
\textbf{\underline{2. Modélisation pour Monsieur Atangana — Objectif de bénéfice (5\,000 FCFA)}}

\noindent \textbf{a) Écriture de l'inéquation.}
Bénéfice $B_A(x) = CA_A(x) - \text{Investissement} = 200x - 8000$.
Objectif : $B_A(x) \geq 5000$ (``au moins'' signifie $\geq$).
\[ 200x - 8000 \geq 5000 \]

\noindent \textbf{b) Rappel définition :} Déjà fait ci-dessus : $B_A(x) = 200x - 8000$.

\noindent \textbf{c) Résolution dans $\mathbb{R}$.}
\begin{align*}
200x - 8000 &\geq 5000 \\
200x - 8000 + 8000 &\geq 5000 + 8000 \qquad \text{(Ajout de 8000 des deux côtés, sens conservé)} \\
200x &\geq 13000 \\
\dfrac{200x}{200} &\geq \dfrac{13000}{200} \qquad \text{(Division par 200 $> 0$, sens \textbf{conservé})} \\
x &\geq 65
\end{align*}
\emph{Attention :} Si on avait divisé par un nombre négatif (ex: $-200$), il aurait fallu inverser le signe $\geq$ en $\leq$. Ici $200 > 0$, pas d'inversion.

\noindent Ensemble solution dans $\mathbb{R}$ : $S_A = [65; +\infty[$.

\noindent \textbf{d) Interprétation dans $\mathbb{N}$.}
Comme $x$ est un nombre entier de brochettes, $x \in \mathbb{N} \cap [65; +\infty[ = \{65, 66, 67, \dots\}$.
Monsieur Atangana doit vendre \textbf{au minimum 65 brochettes} pour gagner au moins 5\,000 FCFA de bénéfice net.

\medskip
\textbf{\underline{3. Comparaison des deux vendeurs — Même affluence ($x$ identique)}}

\noindent \textbf{a) Fonction bénéfice de Monsieur Bello.}
Investissement = 5\,000 FCFA. Prix = 150 FCFA.
\[ B_B(x) = 150x - 5000 \]

\noindent \textbf{b) Hypothèse :} Même nombre $x$ de brochettes vendues par les deux.

\noindent \textbf{c) Inéquation de supériorité stricte.}
On cherche $x$ tel que : Bénéfice Atangana $>$ Bénéfice Bello.
\[ B_A(x) > B_B(x) \]
\[ 200x - 8000 > 150x - 5000 \]

\noindent \textbf{Résolution :}
\begin{align*}
200x - 8000 &> 150x - 5000 \\
200x - 150x - 8000 &> -5000 \qquad \text{(Soustraction de $150x$ des deux côtés, sens conservé)} \\
50x - 8000 &> -5000 \\
50x - 8000 + 8000 &> -5000 + 8000 \qquad \text{(Ajout de 8000 des deux côtés, sens conservé)} \\
50x &> 3000 \\
\dfrac{50x}{50} &> \dfrac{3000}{50} \qquad \text{(Division par 50 $> 0$, sens \textbf{conservé})} \\
x &> 60
\end{align*}

\noindent Ensemble solution dans $\mathbb{R}$ : $S = ]60; +\infty[$.

\noindent \textbf{d) Interprétation économique.}
Dans $\mathbb{N}$, $x \in \{61, 62, 63, \dots\}$.
\textbf{Conclusion :} Si les deux vendeurs vendent \textbf{60 brochettes ou moins}, Monsieur Bello (stratégie bas prix / bas investissement) est plus rentable ou égal.
À partir de \textbf{61 brochettes vendues}, Monsieur Atangana (stratégie haut prix / haut investissement) devient \textbf{strictement plus rentable}.
Le point de bascule (indifférence) est à $x = 60$ brochettes (bénéfices égaux : $200\times60-8000 = 4000$ et $150\times60-5000 = 4000$).

\medskip
\textbf{\underline{4. Synthèse et Conclusion Globale}}

\begin{center}
\begin{tabularx}{\linewidth}{|l|X|X|}\hline
\textbf{Indicateur} & \textbf{Valeur seuil (dans $\mathbb{N}$)} & \textbf{Interprétation pour M. Atangana} \\ \hline
Seuil de rentabilité (Point mort) & $x = 40$ brochettes & En dessous : perte. À 40 : équilibre (0 FCFA gain). \\ \hline
Objectif Bénéfice $\geq$ 5\,000 FCFA & $x \geq 65$ brochettes & Il faut vendre 65 brochettes minimum pour gagner 5\,000 FCFA nets. \\ \hline
Supériorité sur concurrent (M. Bello) & $x \geq 61$ brochettes & Si affluence forte ($\geq 61$), sa stratégie « premium » rapporte plus. \\ \hline
\end{tabularx}
\end{center}

\noindent \textbf{Conclusion rédigée :}
Pour que son commerce de soya à Nkol-Eton soit viable, Monsieur Atangana doit impérativement vendre \textbf{au moins 40 brochettes} par soirée pour ne pas perdre son capital de 8\,000 FCFA. S'il souhaite se verser un salaire net de 5\,000 FCFA (bénéfice), il doit viser \textbf{65 brochettes}. Comparé à son voisin Monsieur Bello qui mise sur le volume à bas prix (150 FCFA), la stratégie d'Atangana (200 FCFA) devient plus avantageuse financièrement dès que l'affluence permet de vendre \textbf{61 brochettes ou plus}. En deçà, la stratégie « petit investissement / petit prix » de Bello est moins risquée et plus rémunératrice. Cette analyse montre l'importance de bien connaître sa clientèle potentielle (volume de vente prévisible) avant de choisir son positionnement prix/qualité.

\medskip
\noindent \textbf{Présentation orale (Points clés pour le rapporteur) :}
\begin{itemize}
    \item « Nous avons posé $x$ = nombre de brochettes. »
    \item « Pour l'inéquation $200x - 8000 \geq 5000$, nous avons ajouté 8000 (règle : ajout même nombre, sens inchangé), puis divisé par 200 (positif, sens inchangé). Si le coefficient de $x$ avait été négatif, nous aurions inversé le signe. »
    \item « Le résultat $x > 60$ signifie qu'au-delà de 60 brochettes, le prix unitaire plus élevé compense l'investissement initial plus lourd. C'est le concept de \emph{levier opérationnel} simplifié. »
\end{itemize}
\end{exoresolu}

\courssub{Bilan de l'activité}
Cette activité d'intégration a permis de mettre en évidence les points suivants, essentiels pour la réussite au BEPC et pour une culture mathématique utile :

\begin{enumerate}
    \item \textbf{La modélisation est une compétence à part entière :} Choisir l'inconnue ($x$ = nombre de brochettes), définir les fonctions $CA(x)$ et $B(x)$, et traduire les phrases « au moins », « strictement supérieur », « couvre exactement » en symboles ($\geq, >, =$) est l'étape la plus délicate et la plus notée à l'examen. Sans traduction correcte, la résolution technique, même parfaite, ne vaut rien.
    \item \textbf{La rigueur des règles de transformation :} L'activité a forcé la distinction entre équation (division par 200, 50 : nombres positifs, pas de piège) et inéquation. Le rappel explicite « division par un positif $\Rightarrow$ sens conservé » dans le corrigé ancre le réflexe qui évite l'erreur fatale (oublier d'inverser le signe si coefficient négatif).
    \item \textbf{L'interprétation dans le bon ensemble ($\mathbb{N}$ vs $\mathbb{R}$) :} La résolution se fait dans $\mathbb{R}$ (continu) pour utiliser les outils algébriques, mais la réponse finale doit revenir dans $\mathbb{N}$ (discret, nombre d'objets). Le passage de $x \geq 65$ à « 65 brochettes minimum » et de $x > 60$ à « 61 brochettes minimum » est une compétence de communication mathématique attendue.
    \item \textbf{Le sens économique des seuils :} Les élèves ont manipulé trois concepts clés de gestion : \emph{Point mort} (équilibre), \emph{Objectif de résultat} (bénéfice cible), \emph{Point d'indifférence / Bascule} (choix entre deux structures de coûts). Cela donne du sens à l'algèbre abstraite.
    \item \textbf{La communication orale et écrite :} La structure imposée (Inconnue $\rightarrow$ Traduction $\rightarrow$ Résolution $\rightarrow$ Interprétation $\rightarrow$ Conclusion) est le standard de rédaction attendu au BEPC (épreuve écrite) et à l'oral. L'exercice de présentation au tableau développe la compétence « Rédiger et justifier une démarche ».
\end{enumerate}

\begin{attention}[Pièges fréquents relevés pendant l'activité]
\begin{itemize}
    \item \textbf{Oublier l'unité :} Écrire $x = 40$ sans dire « 40 brochettes ». Perte de points systématique.
    \item \textbf{Confondre Chiffre d'Affaires et Bénéfice :} Écrire $200x \geq 5000$ pour la question 2 (objectif bénéfice) au lieu de $200x - 8000 \geq 5000$.
    \item \textbf{Gérer le « strictement supérieur » vs « supérieur ou égal » :} Pour la comparaison (Q3), l'inéquation est stricte ($>$), la solution est $]60; +\infty[$, donc le premier entier est 61. Pour le bénéfice (Q2), c'est large ($\geq$), solution $[65; +\infty[$, premier entier 65. La notation d'intervalle et le premier entier admissible doivent être cohérents.
    \item \textbf{Ne pas justifier la conservation du sens :} Écrire « on divise par 200 » sans préciser « 200 est positif, le sens ne change pas » est une rédaction incomplète au BEPC.
\end{itemize}
\end{attention}

\begin{savaistu}[Le saviez-vous ? — Le seuil de rentabilité dans la vraie vie]
Au Cameroun, comme partout, tout entrepreneur (du vendeur de soya au PDG d'une usine d'aluminium à Edéa) calcule son \textbf{Point Mort} (Break-even Point).
La formule générale est : $\text{Quantité}_{\text{seuil}} = \dfrac{\text{Charges Fixes}}{\text{Prix de vente} - \text{Coût Variable Unitaire}}$.
Dans notre exercice simplifié, nous avons supposé \textbf{Coût Variable Unitaire = 0} (pas de coût d'achat de viande par brochette, ni de piment, ni d'huile par unité). C'est pour cela que l'équation était $200x = 8000$.
Dans la réalité, si une brochette coûte 50 FCFA en viande/épices (coût variable), la marge unitaire n'est plus 200 mais $200 - 50 = 150$ FCFA. Le seuil de rentabilité devient $\frac{8000}{150} \approx 53,3 \rightarrow 54$ brochettes. L'algèbre du 1er degré reste exactement la même : $(Prix - CoûtVar) \times x = ChargesFixes$. C'est la puissance de la modélisation mathématique : un même outil (équation 1er degré) gère la soya et l'usine !
\end{savaistu}

\newpage

\coursec{Méthodes et exercices résolus}

\coursec{Équations et inéquations du 1er degré à une inconnue dans $\mathbb{R}$}

\noindent Bienvenue dans cette leçon fondamentale de l'unité d'algèbre. Nous allons consolider la résolution des équations du premier degré, aborder les inéquations — où le signe du coefficient change la donne — et surtout apprendre à modéliser des situations concrètes (achats au marché, gestion de budget, comparaisons d'âges) pour en extraire la réponse mathématique exacte. La rigueur dans l'écriture des étapes et la phrase de conclusion finale sont vos meilleurs atouts pour le BEPC.

\courssub{Résoudre une équation de la forme $ax+b=0$}

\begin{definition}[Équation du premier degré à une inconnue]
Une \cle{équation du premier degré à une inconnue} $x$ s'écrit sous la forme canonique $ax+b=0$, où $a$ et $b$ sont des nombres réels donnés, avec $a \neq 0$. La \cle{solution} (ou racine) est la valeur de $x$ qui rend l'égalité vraie. L'\cle{ensemble des solutions} se note $S$.
\end{definition}

\begin{methode}[Résoudre $ax+b=0$]
Pour résoudre une équation de la forme $ax+b=0$ d'inconnue $x$ :
\begin{enumerate}
\item \textbf{Isoler le terme en $x$} : on ajoute $-b$ (ou on soustrait $b$) aux deux membres : $ax=-b$.
\item \textbf{Diviser par $a$} (rappel : $a\neq 0$) : $x=-\dfrac{b}{a}$.
\item \textbf{Simplifier} la fraction obtenue si possible (recherche de diviseurs communs).
\item \textbf{Vérifier} en remplaçant $x$ par la valeur trouvée dans l'équation de départ.
\item \textbf{Conclure} par une phrase : « L'équation admet pour solution ... » et écrire l'ensemble des solutions $S=\left\{-\dfrac{b}{a}\right\}$.
\end{enumerate}
\textbf{Réflexe :} Toute opération faite sur un membre doit être faite sur l'autre. On ne « déplace » jamais un terme sans changer son signe ; on \emph{ajoute le même nombre aux deux membres}.
\end{methode}

\begin{exemplebox}[Cas particuliers — Rappel]
Si après simplification on obtient :
\begin{itemize}
\item $0x = 0$ (ex: $2x+3 = 2x+3$) : tout réel est solution, $S = \mathbb{R}$ (équation \cle{indéterminée}).
\item $0x = c$ avec $c \neq 0$ (ex: $2x+3 = 2x+5 \Rightarrow 3=5$) : aucune solution, $S = \varnothing$ (équation \cle{impossible}).
\end{itemize}
Ces cas n'apparaissent pas dans la forme $ax+b=0$ avec $a\neq 0$, mais surgissent souvent quand on « ramène » une équation complexe.
\end{exemplebox}

\begin{exoresolu}[Équation du type $ax+b=0$]
Résoudre dans $\mathbb{R}$ l'équation : $3x-12=0$.
\tcblower
\textbf{Étape 1 :} On isole le terme en $x$ en ajoutant $12$ aux deux membres :
\[3x-12+12=0+12 \quad\Longleftrightarrow\quad 3x=12.\]
\textbf{Étape 2 :} On divise les deux membres par $3$ (qui est non nul) :
\[x=\dfrac{12}{3}=4.\]
\textbf{Étape 3 : Vérification.} Pour $x=4$ : $3\times 4-12=12-12=0$. L'égalité est vraie.
\textbf{Conclusion :} L'équation $3x-12=0$ admet une unique solution : $x=4$. On écrit $S=\{4\}$.
\end{exoresolu}

\begin{exoresolu}[Coefficient négatif]
Résoudre dans $\mathbb{R}$ : $-5x+15=0$.
\tcblower
On ajoute $-15$ aux deux membres : $-5x = -15$.
On divise par $-5$ : $x = \dfrac{-15}{-5} = 3$.
\textbf{Vérification :} $-5\times 3 + 15 = -15+15=0$.
\textbf{Conclusion :} $S=\{3\}$.
\end{exoresolu}

\courssub{Résoudre une équation se ramenant au premier degré}

\begin{methode}[Se ramener à la forme $ax+b=0$]
Pour résoudre une équation contenant des $x$ et des nombres dans les deux membres, éventuellement avec des parenthèses ou des dénominateurs :
\begin{enumerate}
\item \textbf{Supprimer les dénominateurs} s'il y en a : multiplier les deux membres par le PPCM des dénominateurs (et noter la condition de non-nullité).
\item \textbf{Développer} et supprimer les parenthèses.
\item \textbf{Réduire} chaque membre (regrouper les termes en $x$ d'un côté, les nombres de l'autre).
\item \textbf{Rassembler les $x$ à gauche, les nombres à droite} en ajoutant/soustrayant les mêmes termes des deux côtés.
\item Se ramener à $ax=b$, puis \textbf{diviser par $a$} ($a\neq 0$).
\item \textbf{Vérifier} que la solution ne rend aucun dénominateur nul (si étape 1 utilisée) et qu'elle satisfait l'équation initiale.
\item \textbf{Conclure} avec $S=\{...\}$.
\end{enumerate}
\textbf{Cas particuliers :} Si on obtient $0x=0 \Rightarrow S=\mathbb{R}$ ; si $0x=c\ (c\neq 0) \Rightarrow S=\varnothing$.
\end{methode}

\begin{exoresolu}[Équation avec des $x$ dans les deux membres]
Résoudre dans $\mathbb{R}$ : $5x-3=2x+9$.
\tcblower
\textbf{Étape 1 :} On soustrait $2x$ aux deux membres (rassembler les $x$ à gauche) :
\[5x-3-2x=2x+9-2x \quad\Longleftrightarrow\quad 3x-3=9.\]
\textbf{Étape 2 :} On ajoute $3$ aux deux membres (rassembler les nombres à droite) :
\[3x-3+3=9+3 \quad\Longleftrightarrow\quad 3x=12.\]
\textbf{Étape 3 :} On divise par $3$ : $x=4$.
\textbf{Vérification :} Gauche : $5\times 4-3=17$ ; Droite : $2\times 4+9=17$. Égalité vraie.
\textbf{Conclusion :} $S=\{4\}$.
\end{exoresolu}

\begin{exoresolu}[Équation avec parenthèses]
Résoudre dans $\mathbb{R}$ : $2(x-3)+5=3(x+1)-8$.
\tcblower
\textbf{Étape 1 : Développement.}
\[2x-6+5=3x+3-8.\]
\textbf{Étape 2 : Réduction de chaque membre.}
\[2x-1=3x-5.\]
\textbf{Étape 3 :} On soustrait $2x$ aux deux membres : $-1=x-5$.
On ajoute $5$ aux deux membres : $4=x$, c'est-à-dire $x=4$.
\textbf{Vérification :} Gauche : $2(4-3)+5=7$ ; Droite : $3(4+1)-8=7$.
\textbf{Conclusion :} $S=\{4\}$.
\end{exoresolu}

\begin{exoresolu}[Équation avec dénominateur — Se ramener au 1er degré]
Résoudre dans $\mathbb{R}$ : $\dfrac{2x+1}{3}=x-2$.
\tcblower
\textbf{Condition :} Le dénominateur $3$ n'est jamais nul, pas de restriction.
\textbf{Étape 1 :} On multiplie les deux membres par $3$ (PPCM) :
\[2x+1 = 3(x-2).\]
\textbf{Étape 2 : Développement.}
\[2x+1 = 3x-6.\]
\textbf{Étape 3 :} On soustrait $2x$ : $1 = x-6$.
On ajoute $6$ : $x = 7$.
\textbf{Vérification :} Gauche : $\frac{2\times 7+1}{3}=\frac{15}{3}=5$ ; Droite : $7-2=5$.
\textbf{Conclusion :} $S=\{7\}$.
\end{exoresolu}

\courssub{Appliquer les règles de transformation des inégalités}

\begin{propriete}[Règles de transformation des inégalités]
Soient $A, B, C \in \mathbb{R}$.
\begin{enumerate}
\item \textbf{Ajouter/Soustraire} un même nombre $C$ : $A < B \Longleftrightarrow A+C < B+C$ (le sens \textbf{ne change pas}).
\item \textbf{Multiplier/Diviser} par un nombre \textbf{strictement positif} $k>0$ : $A < B \Longleftrightarrow kA < kB$ (le sens \textbf{ne change pas}).
\item \textbf{Multiplier/Diviser} par un nombre \textbf{strictement négatif} $k<0$ : $A < B \Longleftrightarrow kA > kB$ (le sens \textbf{change} : $<$ devient $>$, $\leq$ devient $\geq$).
\end{enumerate}
Ces règles valent pour $<, >, \leq, \geq$.
\end{propriete}

\begin{exoresolu}[Transformer une inégalité]
On sait que $x<y$. Comparer, en justifiant : $x-5$ et $y-5$ ; $-3x$ et $-3y$.
\tcblower
\textbf{1.} On soustrait $5$ (règle 1) : $x-5 < y-5$.
\textbf{2.} On multiplie par $-3$ (négatif, règle 3) : le sens change : $-3x > -3y$.
\textbf{Conclusion :} $x-5<y-5$ et $-3x>-3y$.
\end{exoresolu}

\begin{attention}[Piège classique]
Ne jamais écrire « on passe $-4$ de l'autre côté, ça fait $+4$ » dans une inéquation sans préciser le signe. C'est la \textbf{division par $-4$} qui impose le changement de sens. Toujours écrire l'étape : « Je divise par $-4$ (négatif), donc je change le sens. »
\end{attention}

\courssub{Résoudre une inéquation de la forme $ax+b > 0$}

\begin{methode}[Résoudre $ax+b > 0$ (et $<, \leq, \geq$)]
\begin{enumerate}
\item \textbf{Isoler le terme en $x$} : $ax > -b$ (en ajoutant $-b$ aux deux membres, sens inchangé).
\item \textbf{Diviser par $a$} en regardant \textbf{le signe de $a$} :
\begin{itemize}
\item si $a>0$ : le sens ne change pas $\Rightarrow x > -\dfrac{b}{a}$ ;
\item si $a<0$ : le sens \textbf{change} $\Rightarrow x < -\dfrac{b}{a}$.
\end{itemize}
\item \textbf{Écrire l'ensemble des solutions} :
\begin{itemize}
\item Inégalité stricte ($<$ ou $>$) : intervalle ouvert $]-\infty; k[$ ou $]k; +\infty[$, la valeur $k$ est \textbf{exclue}.
\item Inégalité large ($\leq$ ou $\geq$) : intervalle fermé $]-\infty; k]$ ou $[k; +\infty[$, la valeur $k$ est \textbf{incluse}.
\end{itemize}
\item \textbf{Représenter sur une droite graduée} : trait plein pour les solutions, cercle \textbf{vide} (exclu) ou \textbf{plein} (inclus) à la frontière, flèche vers les solutions.
\end{enumerate}
\end{methode}

\begin{popfigure}
\begin{tikzpicture}[scale=0.8, >=Stealth]
% Axe pour x > -3
\draw[popline, thick] (-5,1) -- (3,1);
\foreach \x in {-4,-3,-2,-1,0,1,2}
  \draw (\x,0.9) -- (\x,1.1) node[below=2pt] {\small $\x$};
\draw[popblue, very thick] (-3,1) -- (2.5,1);
\draw[popblue, fill=white, thick] (-3,1) circle (3pt); % cercle vide
\draw[popblue, ->, thick] (2.5,1) -- (2.8,1);
\node[popblue, above] at (-0.5,1.4) {$x > -3 \quad S = ]-3; +\infty[$};

% Axe pour x >= 2
\draw[popline, thick] (-1,-1.5) -- (5,-1.5);
\foreach \x in {0,1,2,3,4}
  \draw (\x,-1.6) -- (\x,-1.4) node[below=2pt] {\small $\x$};
\draw[poporange, very thick] (2,-1.5) -- (4.5,-1.5);
\draw[poporange, fill=poporange, thick] (2,-1.5) circle (3pt); % cercle plein
\draw[poporange, ->, thick] (4.5,-1.5) -- (4.8,-1.5);
\node[poporange, above] at (3,-1.1) {$x \geq 2 \quad S = [2; +\infty[$};
\end{tikzpicture}
\legende{Représentation graphique des ensembles de solutions sur une droite graduée.}
\end{popfigure}

\begin{exoresolu}[Inéquation avec coefficient positif]
Résoudre dans $\mathbb{R}$ : $2x+6>0$.
\tcblower
\textbf{Étape 1 :} On soustrait $6$ : $2x > -6$.
\textbf{Étape 2 :} On divise par $2$ (positif) : $x > -3$.
\textbf{Conclusion :} $S = \left]-3\,;\,+\infty\right[$. La valeur $-3$ n'est \textbf{pas} solution (inégalité stricte, cercle vide sur la droite).
\end{exoresolu}

\begin{exoresolu}[Inéquation avec coefficient négatif]
Résoudre dans $\mathbb{R}$ : $-4x+8\leq 0$.
\tcblower
\textbf{Étape 1 :} On soustrait $8$ : $-4x \leq -8$.
\textbf{Étape 2 :} On divise par $-4$ (négatif) $\Rightarrow$ \textbf{changement de sens} : $x \geq 2$.
\textbf{Vérification :} $x=2 \Rightarrow -8+8=0 \leq 0$ (vrai, inclus). $x=3 \Rightarrow -12+8=-4 \leq 0$ (vrai). $x=1 \Rightarrow -4+8=4 > 0$ (faux, $1$ n'est pas solution).
\textbf{Conclusion :} $S = \left[2\,;\,+\infty\right[$. La valeur $2$ est solution (inégalité large, cercle plein).
\end{exoresolu}

\courssub{Résoudre un problème concret par une équation ou une inéquation}

\begin{methode}[Modéliser et résoudre un problème]
\begin{enumerate}
\item \textbf{Choisir l'inconnue} : « Soit $x$ ... » (préciser l'unité : nombre d'objets, prix en FCFA, âge en années...).
\item \textbf{Traduire} les contraintes en équation ($=$) ou inéquation ($<, \leq, >, \geq$).
\item \textbf{Résoudre} l'équation ou l'inéquation (méthodes ci-dessus).
\item \textbf{Interpréter} le résultat dans le contexte : $x$ doit-il être entier ? positif ? La solution est-elle unique ou un intervalle ?
\item \textbf{Conclure} par une phrase complète répondant à la question, avec l'unité.
\end{enumerate}
\end{methode}

\begin{exoresolu}[Problème « Budget marché » — Inéquation]
Mama Fotso achète des tomates au marché de Mokolo. Elle paie un transport fixe de $500$ FCFA, puis achète des barquettes à $250$ FCFA l'unité. Elle dispose de $3\,000$ FCFA au maximum. Combien de barquettes peut-elle acheter au plus ?
\tcblower
\textbf{1. Inconnue :} Soit $x$ le nombre de barquettes achetées ($x \in \mathbb{N}$).
\textbf{2. Inéquation :} Dépense totale $\leq$ Budget $\Rightarrow 500 + 250x \leq 3\,000$.
\textbf{3. Résolution :}
$250x \leq 2\,500$ (soustraction de $500$)
$x \leq 10$ (division par $250>0$, sens inchangé).
\textbf{4. Interprétation :} $x$ entier naturel $\leq 10$. Le maximum est $10$.
\textbf{Conclusion :} Mama Fotso peut acheter \textbf{au plus 10 barquettes} de tomates.
\end{exoresolu}

\begin{exoresolu}[Problème « Âges » — Équation]
Un père a $40$ ans et son fils a $10$ ans. Dans combien d'années l'âge du père sera-t-il le double de l'âge du fils ?
\tcblower
\textbf{1. Inconnue :} Soit $x$ le nombre d'années cherché ($x \in \mathbb{N}$).
\textbf{2. Équation :} Dans $x$ ans, père : $40+x$, fils : $10+x$. Condition : $40+x = 2(10+x)$.
\textbf{3. Résolution :}
$40+x = 20+2x$
$40-20 = 2x-x$ (on rassemble $x$ à droite, nombres à gauche)
$20 = x$.
\textbf{4. Vérification :} Dans $20$ ans, père $60$ ans, fils $30$ ans. $60 = 2 \times 30$. Exact.
\textbf{Conclusion :} Dans \textbf{20 ans}, l'âge du père sera le double de celui de son fils.
\end{exoresolu}

\begin{exoresolu}[Problème « Forfaits téléphoniques » — Choix entre deux offres]
Orange propose un forfait à $2\,000$ FCFA/mois + $50$ FCFA/min d'appel. MTN propose $3\,500$ FCFA/mois + $30$ FCFA/min. À partir de combien de minutes MTN devient-il moins cher ?
\tcblower
\textbf{1. Inconnue :} Soit $x$ le nombre de minutes d'appel par mois ($x \geq 0$).
\textbf{2. Inéquation :} Coût MTN $<$ Coût Orange $\Rightarrow 3\,500 + 30x < 2\,000 + 50x$.
\textbf{3. Résolution :}
$3\,500 - 2\,000 < 50x - 30x$
$1\,500 < 20x$
Division par $20>0$ : $75 < x$ c'est-à-dire $x > 75$.
\textbf{4. Interprétation :} Pour $x=75$, coûts égaux ($3\,500+2\,250=5\,750$). Pour $x=76$, MTN moins cher.
\textbf{Conclusion :} MTN devient moins cher \textbf{à partir de 76 minutes} d'appel par mois.
\end{exoresolu}

\begin{savaistu}[Al-Khwârizmî, le père de l'algèbre]
Le mot « algèbre » vient de l'arabe \emph{al-jabr} (restauration, complétion), titre de l'ouvrage \emph{Al-Kitāb al-mukhtaṣar fī ḥisāb al-jabr wal-muqābala} (vers 820) du mathématicien perse \textbf{Muhammad ibn Mūsā al-Khwârizmî}. Il y décrivait les méthodes systématiques pour résoudre les équations du 1er et 2nd degré par « réduction » (al-jabr) et « équilibrage » (al-muqābala) — exactement ce que nous faisons en ajoutant le même terme des deux côtés ! L'algorithme (son nom latinisé) et l'algèbre sont nés là. Au Cameroun, ces outils modélisent le calcul des coûts de transport (CAMRAIL), la facturation d'Eneo ou la gestion de crédit MTN/Orange Money.
\end{savaistu}

\begin{attention}[Les 5 erreurs qui coûtent des points au BEPC]
\begin{enumerate}
\item \textbf{Oublier de changer le sens de l'inégalité} en divisant par un nombre négatif : de $-4x\leq -8$, on tire $x\geq 2$ et \textbf{pas} $x\leq 2$. Toujours vérifier le signe du coefficient \textbf{avant} de diviser.
\item \textbf{« Déplacer » un terme sans changer son signe} : de $3x-12=0$, écrire $3x=-12$ est faux ; on obtient $3x=12$. Rédiger : « On ajoute $12$ aux deux membres ».
\item \textbf{Inclure la valeur frontière dans une inégalité stricte} : pour $x>-3$, la valeur $-3$ n'est \textbf{pas} solution. Ne pas confondre $>$ et $\geq$, ni les crochets $]$ et $[$ dans les intervalles.
\item \textbf{Ne pas vérifier la solution} : remplacer $x$ par la valeur trouvée dans l'équation de départ élimine les erreurs de calcul ; c'est un réflexe attendu au BEPC.
\item \textbf{Oublier la phrase de conclusion et l'unité} dans les problèmes : « $x=10$ » ne suffit pas ; il faut écrire « Mama Fotso peut acheter au plus 10 barquettes ». Une réponse sans phrase ni unité perd des points de présentation et de rigueur.
\end{enumerate}
\end{attention}

\begin{exercice}[Entraînement — Équations]
Résoudre dans $\mathbb{R}$ les équations suivantes :
\begin{enumerate}
\item $7x - 21 = 0$
\item $4x + 5 = 2x - 7$
\item $3(x-2) - 2(x+1) = 5$
\item $\dfrac{x+2}{4} = \dfrac{2x-1}{3}$
\end{enumerate}
\end{exercice}

\begin{corrige}[Exercice Équations]
\begin{enumerate}
\item $7x=21 \Rightarrow x=3$. $S=\{3\}$.
\item $4x-2x = -7-5 \Rightarrow 2x=-12 \Rightarrow x=-6$. $S=\{-6\}$.
\item $3x-6-2x-2=5 \Rightarrow x-8=5 \Rightarrow x=13$. $S=\{13\}$.
\item $3(x+2)=4(2x-1) \Rightarrow 3x+6=8x-4 \Rightarrow 10=5x \Rightarrow x=2$. $S=\{2\}$.
\end{enumerate}
\end{corrige}

\begin{exercice}[Entraînement — Inéquations]
Résoudre dans $\mathbb{R}$ et représenter les solutions sur une droite graduée :
\begin{enumerate}
\item $-3x + 9 > 0$
\item $5x - 10 \leq 0$
\item $2(x-4) < 3x + 1$
\end{enumerate}
\end{exercice}

\begin{corrige}[Exercice Inéquations]
\begin{enumerate}
\item $-3x > -9 \Rightarrow$ division par $-3$ (négatif) $\Rightarrow x < 3$. $S=]-\infty; 3[$.
\item $5x \leq 10 \Rightarrow$ division par $5$ (positif) $\Rightarrow x \leq 2$. $S=]-\infty; 2]$.
\item $2x-8 < 3x+1 \Rightarrow -x < 9 \Rightarrow$ division par $-1$ $\Rightarrow x > -9$. $S=]-9; +\infty[$.
\end{enumerate}
\end{corrige}

\begin{exercice}[Problème de synthèse]
Un car de transport « Voyage Confort » part de Yaoundé pour Douala. Le prix du billet est de $4\,500$ FCFA. Le carburant et les péages coûtent $150\,000$ FCFA fixes pour le trajet. Le chauffeur veut faire un bénéfice d'au moins $50\,000$ FCFA. Combien de passagers minimum doit-il transporter ? (On suppose le car plein ou non, pas de limite de places).
\end{exercice}

\begin{corrige}[Problème de synthèse]
\textbf{Inconnue :} $x$ nombre de passagers ($x \in \mathbb{N}^*$).
\textbf{Inéquation :} Recettes $-$ Dépenses $\geq$ Bénéfice visé.
$4\,500x - 150\,000 \geq 50\,000$.
$4\,500x \geq 200\,000$.
$x \geq \dfrac{200\,000}{4\,500} = \dfrac{400}{9} \approx 44,44$.
\textbf{Interprétation :} $x$ entier $\geq 44,44 \Rightarrow x \geq 45$.
\textbf{Conclusion :} Il doit transporter \textbf{au minimum 45 passagers}.
\end{corrige}

\newpage

\coursec{Exercices}

\courssub{Je vérifie mes connaissances}

\begin{exercice}[Résolution directe d'équations]
Résous dans $\mathbb{R}$ les équations suivantes :
\begin{enumerate}
    \item $4x - 12 = 0$
    \item $7x + 3 = 2x - 17$
    \item $5(x - 2) = 3(x + 4)$
    \item $\dfrac{x}{2} + 1 = \dfrac{x}{3} + 2$
\end{enumerate}
\end{exercice}

\begin{exercice}[Vrai ou Faux ? Justifie]
Pour chaque affirmation, dis si elle est \textbf{Vraie} ou \textbf{Fausse} et justifie ta réponse par un calcul.
\begin{enumerate}
    \item $-3$ est solution de l'équation $2x + 6 = 0$.
    \item $5$ est solution de l'inéquation $-2x + 8 > 0$.
    \item L'équation $3x - 9 = 0$ a pour solution $x = -3$.
    \item L'inéquation $4x - 12 \leqslant 0$ a pour ensemble de solutions $]-\infty\,; 3]$.
\end{enumerate}
\end{exercice}

\begin{exercice}[Vocabulaire et définitions]
Complète les phrases suivantes en utilisant les termes appropriés : \emph{solution, inconnue, équivalent, ensemble solution, vérifier}.
\begin{enumerate}
    \item Dans l'équation $ax + b = 0$, la lettre $x$ s'appelle l'\underline{\hspace{3cm}}.
    \item Une valeur de $x$ qui rend l'égalité vraie s'appelle une \underline{\hspace{3cm}} de l'équation.
    \item Deux équations qui ont le même ensemble de solutions sont dites \underline{\hspace{3cm}}.
    \item L'\underline{\hspace{3cm}} d'une inéquation est l'ensemble de toutes les valeurs qui la vérifient.
\end{enumerate}
\end{exercice}

\begin{exercice}[Résolution d'inéquations simples]
Résous dans $\mathbb{R}$ les inéquations suivantes :
\begin{enumerate}
    \item $5x - 15 > 0$
    \item $-4x + 12 \geqslant 0$
    \item $3x + 7 < 1$
    \item $-2(x - 5) \geqslant 6$
\end{enumerate}
\end{exercice}

\courssub{Je m'entraîne}

\begin{exercice}[Équations se ramenant au premier degré]
Résous dans $\mathbb{R}$ les équations suivantes en détaillant les étapes (développement, regroupement, division) :
\begin{enumerate}
    \item $(3x - 9)(2x + 5) = 0$ \quad (Attention : produit nul)
    \item $\dfrac{x + 1}{3} = \dfrac{2x - 5}{2}$
    \item $2(3x - 4) - 5(x + 1) = 3$
    \item $\dfrac{2x - 1}{4} - \dfrac{x + 2}{3} = 1$
\end{enumerate}
\end{exercice}

\begin{exercice}[Inéquations et représentation graphique]
Résous les inéquations suivantes dans $\mathbb{R}$ et représente l'ensemble des solutions sur une droite graduée.
\begin{enumerate}
    \item $2x - 8 \leqslant 4$
    \item $-3x + 9 > 0$
    \item $5 - 2x \geqslant 11$
    \item $\dfrac{x}{2} + 3 < \dfrac{x}{3} + 5$
\end{enumerate}
\end{exercice}

\begin{exercice}[Problème concret : Achat de carnets]
Dans une librairie à Yaoundé, un cahier de 100 pages coûte 350 FCFA et un cahier de 200 pages coûte 600 FCFA. Un élève achète au total 8 cahiers pour une somme de 3 800 FCFA.
\begin{enumerate}
    \item Soit $x$ le nombre de cahiers de 100 pages achetés. Exprime en fonction de $x$ le nombre de cahiers de 200 pages.
    \item Écris l'équation traduisit le prix total.
    \item Résous l'équation et déduis le nombre de cahiers de chaque type achetés.
\end{enumerate}
\end{exercice}

\begin{exercice}[Comparaison de deux offres : Forfaits Internet]
Deux opérateurs proposent des forfaits mensuels pour l'internet mobile :
\begin{itemize}
    \item \textbf{Offre MTN :} 2 000 FCFA par mois + 50 FCFA par Mo consommé au-delà du forfait de base.
    \item \textbf{Offre Orange :} 3 500 FCFA par mois (illimité).
\end{itemize}
Soit $x$ le nombre de Mo consommés en plus du forfait de base chez MTN.
\begin{enumerate}
    \item Écris le prix $P_{\text{MTN}}(x)$ à payer chez MTN en fonction de $x$.
    \item Écris le prix $P_{\text{Orange}}(x)$ à payer chez Orange en fonction de $x$.
    \item Pour quelle consommation $x$ les deux offres coûtent-elles le même prix ?
    \item Détermine l'ensemble des valeurs de $x$ pour lesquelles l'offre MTN est moins chère que l'offre Orange.
\end{enumerate}
\end{exercice}

\courssub{Je me prépare au Probatoire}

\begin{exercice}[Taxi-moto : Gestion d'un budget]
Un taxi-moto facture une prise en charge de 300 FCFA plus 100 FCFA par kilomètre parcouru. Un client dispose de 2 000 FCFA pour son trajet.
\begin{enumerate}
    \item Soit $x$ la distance parcourue en kilomètres. Exprime le prix $P(x)$ de la course en fonction de $x$.
    \item Le client ne doit pas dépasser son budget. Écris l'inéquation que doit vérifier $x$.
    \item Résous cette inéquation dans $\mathbb{R}$.
    \item Sachant que la distance est un nombre positif, quelle est la \textbf{distance maximale} (entière) que le client peut parcourir ? Justifie.
    \item Si le client parcourt 15 km, combien doit-il payer ? Peut-il payer avec ses 2 000 FCFA ?
\end{enumerate}
\end{exercice}

\begin{exercice}[Cybercafés : Choix du meilleur tarif]
Deux cybercafés du quartier proposent les tarifs horaires suivants :
\begin{itemize}
    \item \textbf{Cyber A :} 200 FCFA de droit d'entrée + 150 FCFA par heure de connexion.
    \item \textbf{Cyber B :} 500 FCFA de droit d'entrée + 100 FCFA par heure de connexion.
\end{itemize}
Soit $h$ la durée de connexion en heures ($h > 0$).
\begin{enumerate}
    \item Écris les fonctions $P_A(h)$ et $P_B(h)$ donnant le prix à payer dans chaque cybercafé.
    \item Résous l'équation $P_A(h) = P_B(h)$. Interprète le résultat obtenu.
    \item Résous l'inéquation $P_A(h) < P_B(h)$.
    \item À partir de quelle durée de connexion (en heures et minutes) le Cyber B devient-il plus avantageux que le Cyber A ?
    \item Un élève a 1 500 FCFA. Dans quel(s) cybercafé(s) peut-il se connecter pendant 5 heures ? Justifie par le calcul.
\end{enumerate}
\end{exercice}

\begin{exercice}[Âges : Problème ouvert]
Cette année, un père a 42 ans et son fils a 12 ans.
\begin{enumerate}
    \item Vérifie que le père a actuellement trois fois et demi l'âge de son fils.
    \item Dans $x$ années, quel sera l'âge du père ? Quel sera l'âge du fils ?
    \item On cherche à savoir \textbf{dans combien d'années} l'âge du père sera \textbf{strictement inférieur au double} de l'âge de son fils.
    \begin{enumerate}
        \item Écris l'inéquation traduisant cette situation.
        \item Résous cette inéquation dans $\mathbb{R}$.
        \item Comme $x$ représente un nombre d'années, que peux-tu dire sur les valeurs entières possibles de $x$ ?
        \item Donne la réponse finale en phrase complète.
    \end{enumerate}
    \item Dans combien d'années l'âge du père sera-t-il exactement le double de l'âge de son fils ? (Résous l'équation correspondante).
\end{enumerate}
\end{exercice}

\newpage

\coursec{Corrigés des exercices}

\begin{corrige}[Exercice 1 — Résolution directe d'équations]
\begin{enumerate}
    \item $4x - 12 = 0 \iff 4x = 12 \iff x = \dfrac{12}{4} = 3$. \quad $S = \{3\}$.
    \item $7x + 3 = 2x - 17 \iff 7x - 2x = -17 - 3 \iff 5x = -20 \iff x = -4$. \quad $S = \{-4\}$.
    \item $5(x - 2) = 3(x + 4) \iff 5x - 10 = 3x + 12 \iff 5x - 3x = 12 + 10 \iff 2x = 22 \iff x = 11$. \quad $S = \{11\}$.
    \item $\dfrac{x}{2} + 1 = \dfrac{x}{3} + 2$. On multiplie les deux membres par $6$ (PPCM de 2 et 3) :
    \[ 3x + 6 = 2x + 12 \iff 3x - 2x = 12 - 6 \iff x = 6. \]
    $S = \{6\}$.
\end{enumerate}
\end{corrige}

\begin{corrige}[Exercice 2 — Vrai ou Faux ? Justifie]
\begin{enumerate}
    \item \textbf{Vrai.} On remplace $x$ par $-3$ : $2\times(-3) + 6 = -6 + 6 = 0$. L'égalité est vérifiée, donc $-3$ est solution.
    \item \textbf{Faux.} On remplace $x$ par $5$ : $-2\times 5 + 8 = -10 + 8 = -2$, et $-2 > 0$ est faux. Donc $5$ n'est pas solution.
    \item \textbf{Faux.} $3x - 9 = 0 \iff 3x = 9 \iff x = 3$. La solution est $3$ et non $-3$.
    \item \textbf{Vrai.} $4x - 12 \leqslant 0 \iff 4x \leqslant 12 \iff x \leqslant 3$ (on divise par $4 > 0$, le sens ne change pas). L'ensemble des solutions est bien $]-\infty\,; 3]$.
\end{enumerate}
\end{corrige}

\begin{corrige}[Exercice 3 — Vocabulaire et définitions]
\begin{enumerate}
    \item Dans l'équation $ax + b = 0$, la lettre $x$ s'appelle l'\textbf{inconnue}.
    \item Une valeur de $x$ qui rend l'égalité vraie s'appelle une \textbf{solution} de l'équation.
    \item Deux équations qui ont le même ensemble de solutions sont dites \textbf{équivalentes}.
    \item L'\textbf{ensemble solution} d'une inéquation est l'ensemble de toutes les valeurs qui la vérifient.
\end{enumerate}
\end{corrige}

\begin{corrige}[Exercice 4 — Résolution d'inéquations simples]
\begin{enumerate}
    \item $5x - 15 > 0 \iff 5x > 15 \iff x > 3$. \quad $S = ]3\,; +\infty[$.
    \item $-4x + 12 \geqslant 0 \iff -4x \geqslant -12$. On divise par $-4 < 0$ : \textbf{on change le sens} de l'inégalité :
    \[ x \leqslant \dfrac{-12}{-4} = 3. \qquad S = ]-\infty\,; 3]. \]
    \item $3x + 7 < 1 \iff 3x < -6 \iff x < -2$. \quad $S = ]-\infty\,; -2[$.
    \item $-2(x - 5) \geqslant 6 \iff -2x + 10 \geqslant 6 \iff -2x \geqslant -4 \iff x \leqslant 2$ (division par $-2 < 0$, sens changé). \quad $S = ]-\infty\,; 2]$.
\end{enumerate}
\end{corrige}

\begin{corrige}[Exercice 5 — Équations se ramenant au premier degré]
\begin{enumerate}
    \item $(3x - 9)(2x + 5) = 0$. Un produit de facteurs est nul si et seulement si l'un des facteurs est nul :
    \[ 3x - 9 = 0 \iff x = 3 \qquad \text{ou} \qquad 2x + 5 = 0 \iff x = -\dfrac{5}{2}. \]
    $S = \left\{-\dfrac{5}{2}\,; 3\right\}$.
    \item $\dfrac{x + 1}{3} = \dfrac{2x - 5}{2}$. On multiplie les deux membres par $6$ :
    \[ 2(x + 1) = 3(2x - 5) \iff 2x + 2 = 6x - 15 \iff 2 + 15 = 6x - 2x \iff 17 = 4x \iff x = \dfrac{17}{4}. \]
    $S = \left\{\dfrac{17}{4}\right\}$.
    \item $2(3x - 4) - 5(x + 1) = 3 \iff 6x - 8 - 5x - 5 = 3 \iff x - 13 = 3 \iff x = 16$. \quad $S = \{16\}$.
    \item $\dfrac{2x - 1}{4} - \dfrac{x + 2}{3} = 1$. On multiplie par $12$ (PPCM de 4 et 3) :
    \[ 3(2x - 1) - 4(x + 2) = 12 \iff 6x - 3 - 4x - 8 = 12 \iff 2x - 11 = 12 \iff 2x = 23 \iff x = \dfrac{23}{2}. \]
    $S = \left\{\dfrac{23}{2}\right\}$.
\end{enumerate}
\end{corrige}

\begin{corrige}[Exercice 6 — Inéquations et représentation graphique]
\begin{enumerate}
    \item $2x - 8 \leqslant 4 \iff 2x \leqslant 12 \iff x \leqslant 6$. \quad $S = ]-\infty\,; 6]$.
    \item $-3x + 9 > 0 \iff -3x > -9 \iff x < 3$ (division par $-3 < 0$ : sens changé). \quad $S = ]-\infty\,; 3[$.
    \item $5 - 2x \geqslant 11 \iff -2x \geqslant 6 \iff x \leqslant -3$ (sens changé). \quad $S = ]-\infty\,; -3]$.
    \item $\dfrac{x}{2} + 3 < \dfrac{x}{3} + 5$. On multiplie par $6$ : $3x + 18 < 2x + 30 \iff x < 12$. \quad $S = ]-\infty\,; 12[$.
\end{enumerate}

\begin{popfigure}
\begin{tikzpicture}[scale=0.9]
% 1) x <= 6
\draw[->,popline] (-1,0) -- (8,0);
\foreach \x in {0,2,4,6} \draw (\x,0.08) -- (\x,-0.08) node[below]{\small $\x$};
\draw[line width=2pt,popblue] (-0.8,0) -- (6,0);
\draw[popblue,fill=popblue] (6,0) circle (0.09);
\node[popblue,above] at (3,0.1) {\small $x \leqslant 6$};
% 2) x < 3
\begin{scope}[yshift=-1.3cm]
\draw[->,popline] (-1,0) -- (8,0);
\foreach \x in {0,2,4,6} \draw (\x,0.08) -- (\x,-0.08) node[below]{\small $\x$};
\draw[line width=2pt,poporange] (-0.8,0) -- (3,0);
\draw[poporange,fill=white] (3,0) circle (0.09);
\node[poporange,above] at (1,0.1) {\small $x < 3$};
\end{scope}
% 3) x <= -3
\begin{scope}[yshift=-2.6cm]
\draw[->,popline] (-6,0) -- (2,0);
\foreach \x in {-4,-2,0} \draw (\x,0.08) -- (\x,-0.08) node[below]{\small $\x$};
\draw[line width=2pt,popgreen] (-5.8,0) -- (-3,0);
\draw[popgreen,fill=popgreen] (-3,0) circle (0.09);
\node[popgreen,above] at (-4.5,0.1) {\small $x \leqslant -3$};
\end{scope}
% 4) x < 12
\begin{scope}[yshift=-3.9cm]
\draw[->,popline] (-1,0) -- (14,0);
\foreach \x in {0,4,8,12} \draw (\x,0.08) -- (\x,-0.08) node[below]{\small $\x$};
\draw[line width=2pt,poppurple] (-0.8,0) -- (12,0);
\draw[poppurple,fill=white] (12,0) circle (0.09);
\node[poppurple,above] at (5,0.1) {\small $x < 12$};
\end{scope}
\end{tikzpicture}
\legende{Ensembles de solutions : point plein = valeur incluse, point creux = valeur exclue.}
\end{popfigure}
\end{corrige}

\begin{corrige}[Exercice 7 — Problème concret : Achat de carnets]
\begin{enumerate}
    \item L'élève achète 8 cahiers au total, donc le nombre de cahiers de 200 pages est $8 - x$.
    \item Prix total : $350x + 600(8 - x) = 3\,800$.
    \item Résolution :
    \begin{align*}
    350x + 600(8 - x) &= 3\,800 \\
    350x + 4\,800 - 600x &= 3\,800 \\
    -250x &= 3\,800 - 4\,800 \\
    -250x &= -1\,000 \\
    x &= \dfrac{-1\,000}{-250} = 4.
    \end{align*}
    Donc $8 - x = 4$. \textbf{Conclusion :} l'élève a acheté \textbf{4 cahiers de 100 pages} et \textbf{4 cahiers de 200 pages}.\\
    \emph{Vérification :} $4 \times 350 + 4 \times 600 = 1\,400 + 2\,400 = 3\,800$ FCFA. \checkmark
\end{enumerate}
\end{corrige}

\begin{corrige}[Exercice 8 — Comparaison de deux offres : Forfaits Internet]
\begin{enumerate}
    \item $P_{\text{MTN}}(x) = 2\,000 + 50x$.
    \item $P_{\text{Orange}}(x) = 3\,500$ (prix constant, indépendant de $x$).
    \item On résout $P_{\text{MTN}}(x) = P_{\text{Orange}}(x)$ :
    \[ 2\,000 + 50x = 3\,500 \iff 50x = 1\,500 \iff x = 30. \]
    Les deux offres coûtent le même prix pour une consommation supplémentaire de \textbf{30 Mo}.
    \item On résout $P_{\text{MTN}}(x) < P_{\text{Orange}}(x)$ :
    \[ 2\,000 + 50x < 3\,500 \iff 50x < 1\,500 \iff x < 30. \]
    Comme $x \geqslant 0$, l'ensemble des valeurs est $S = [0\,; 30[$.\\
    \textbf{Conclusion :} l'offre MTN est moins chère tant que la consommation supplémentaire est \textbf{strictement inférieure à 30 Mo} ; au-delà, Orange est plus avantageuse.
\end{enumerate}
\end{corrige}

\begin{corrige}[Exercice 9 — Taxi-moto : Gestion d'un budget]
\begin{enumerate}
    \item $P(x) = 300 + 100x$ (en FCFA).
    \item Le budget ne doit pas être dépassé : $300 + 100x \leqslant 2\,000$.
    \item Résolution :
    \[ 100x \leqslant 2\,000 - 300 \iff 100x \leqslant 1\,700 \iff x \leqslant 17. \]
    $S = ]-\infty\,; 17]$.
    \item Comme $x$ est une distance, $x \geqslant 0$, donc $x \in [0\,; 17]$. La distance maximale entière est $\boxed{17 \text{ km}}$.\\
    \emph{Vérification :} $P(17) = 300 + 1\,700 = 2\,000$ FCFA, exactement le budget.
    \item $P(15) = 300 + 100 \times 15 = 300 + 1\,500 = 1\,800$ FCFA. Comme $1\,800 < 2\,000$, le client \textbf{peut payer} avec ses 2 000 FCFA (il lui restera 200 FCFA).
\end{enumerate}

\textbf{Barème indicatif (sur 5 pts) :} expression de $P(x)$ : 1 pt ; inéquation : 1 pt ; résolution avec justification : 1 pt ; distance maximale justifiée ($x \geqslant 0$ cité) : 1 pt ; question 5 avec conclusion : 1 pt.

\textbf{Points de vigilance :} citer la condition $x \geqslant 0$ (une distance est positive) ; ne pas oublier la prise en charge de 300 FCFA dans $P(x)$ ; conclure par une phrase avec l'unité (km, FCFA).
\end{corrige}

\begin{corrige}[Exercice 10 — Cybercafés : Choix du meilleur tarif]
\begin{enumerate}
    \item $P_A(h) = 200 + 150h$ \quad et \quad $P_B(h) = 500 + 100h$ (en FCFA).
    \item $P_A(h) = P_B(h) \iff 200 + 150h = 500 + 100h \iff 50h = 300 \iff h = 6$.\\
    \textbf{Interprétation :} pour \textbf{6 heures} de connexion, les deux cybercafés coûtent le même prix, soit $200 + 150 \times 6 = 1\,100$ FCFA.
    \item $P_A(h) < P_B(h) \iff 200 + 150h < 500 + 100h \iff 50h < 300 \iff h < 6$.\\
    $S = ]0\,; 6[$ (car $h > 0$). Le Cyber A est moins cher pour moins de 6 heures.
    \item Le Cyber B devient plus avantageux lorsque $h > 6$, c'est-à-dire \textbf{à partir de 6 heures et 1 minute de connexion} (strictement au-delà de 6 h).
    \item Pour $h = 5$ :
    \[ P_A(5) = 200 + 150 \times 5 = 950 \text{ FCFA} \qquad ; \qquad P_B(5) = 500 + 100 \times 5 = 1\,000 \text{ FCFA}. \]
    Comme $950 \leqslant 1\,500$ et $1\,000 \leqslant 1\,500$, l'élève peut se connecter 5 heures \textbf{dans les deux cybercafés} (le Cyber A étant le moins cher des deux pour cette durée).
\end{enumerate}

\textbf{Barème indicatif (sur 6 pts) :} fonctions : 1 pt ; équation + interprétation : 1,5 pt ; inéquation : 1 pt ; conversion heures/minutes : 1 pt ; question 5 avec calculs et conclusion : 1,5 pt.

\textbf{Points de vigilance :} bien interpréter le résultat de l'équation en phrase ; respecter la condition $h > 0$ ; pour la question 4, remarquer que l'inégalité est stricte ($h > 6$) : à 6 h exactement, les deux tarifs sont égaux.
\end{corrige}

\begin{corrige}[Exercice 11 — Âges : Problème ouvert]
\begin{enumerate}
    \item $12 \times \num{3,5} = 42$. Le père a bien actuellement trois fois et demi l'âge de son fils.
    \item Dans $x$ années : le père aura $42 + x$ ans et le fils aura $12 + x$ ans.
    \item
    \begin{enumerate}
        \item L'inéquation est : $42 + x < 2(12 + x)$.
        \item Résolution :
        \begin{align*}
        42 + x &< 24 + 2x \\
        42 - 24 &< 2x - x \\
        18 &< x \quad \text{c'est-à-dire} \quad x > 18.
        \end{align*}
        $S = ]18\,; +\infty[$.
        \item Comme $x$ est un nombre d'années, on retient les valeurs \textbf{entières strictement supérieures à 18}, c'est-à-dire $x \geqslant 19$.
        \item \textbf{Réponse :} c'est à partir de \textbf{19 années} (donc dans 19 ans et au-delà) que l'âge du père sera strictement inférieur au double de l'âge de son fils. \emph{Vérification pour $x = 19$ : père : $61$ ans ; fils : $31$ ans ; $2 \times 31 = 62 > 61$.} \checkmark
    \end{enumerate}
    \item On résout l'équation $42 + x = 2(12 + x)$ :
    \[ 42 + x = 24 + 2x \iff 18 = x. \]
    \textbf{Dans 18 ans exactement}, l'âge du père ($60$ ans) sera le double de celui de son fils ($30$ ans). Cela confirme la question 3 : l'inégalité stricte n'est vérifiée qu'après 18 ans.
\end{enumerate}

\textbf{Barème indicatif (sur 5 pts) :} vérification : 0,5 pt ; âges en fonction de $x$ : 1 pt ; inéquation : 0,5 pt ; résolution : 1 pt ; conclusion entière + phrase : 1 pt ; équation de la question 4 : 1 pt.

\textbf{Points de vigilance :} distinguer l'inégalité \emph{stricte} ($<$) de l'égalité : la valeur $x = 18$ est exclue de la question 3 mais est la réponse de la question 4 ; exiger une phrase de conclusion avec le contexte (années, âges).
\end{corrige}

\newpage

\coursec{Fiche bilan}

\begin{popfigure}
\begin{tikzpicture}[mindmap, grow cyclic, every node/.style={concept, font=\small},
root concept/.append style={concept color=popblue, font=\bfseries\small, minimum size=2.6cm},
level 1/.append style={level distance=3.4cm, sibling angle=60, minimum size=1.9cm, font=\scriptsize}]
\node[root concept]{Équations et inéquations du 1\textsuperscript{er} degré}
child[concept color=poporange]{ node[align=center]{Équation $ax+b=0$\\ $a\neq 0$ : $x=-\dfrac{b}{a}$} }
child[concept color=popgreen]{ node[align=center]{Équations s'y ramenant\\ produit nul, fractions} }
child[concept color=poppurple]{ node[align=center]{Règles sur les inégalités\\ sens conservé ou changé} }
child[concept color=poppink]{ node[align=center]{Inéquation $ax+b>0$\\ solutions = intervalle} }
child[concept color=popgold]{ node[align=center]{Tableau de signes\\ de $ax+b$} }
child[concept color=popmuted]{ node[align=center]{Problèmes concrets\\ mise en équation} };
\end{tikzpicture}
\legende{Carte mentale de la leçon : équations et inéquations du premier degré à une inconnue dans $\mathbb{R}$.}
\end{popfigure}

\begin{bilan}
\textbf{Définitions clés}
\begin{itemize}
\item Une \cle{équation du premier degré} à une inconnue $x$ est une égalité pouvant s'écrire sous la forme $ax+b=0$ avec $a\neq 0$.
\item \cle{Résoudre} une équation (ou une inéquation) dans $\mathbb{R}$, c'est trouver \textbf{toutes} les valeurs de $x$ qui la rendent vraie : c'est l'\cle{ensemble des solutions}, noté $\mathcal{S}$.
\item Une \cle{inéquation du premier degré} est une inégalité pouvant s'écrire $ax+b>0$ (ou $<0$, $\geq 0$, $\leq 0$).
\item Un \cle{produit est nul} si et seulement si l'un au moins de ses facteurs est nul : $A\times B=0 \iff A=0$ ou $B=0$.
\end{itemize}

\textbf{Formules et résultats à connaître par cœur}
\begin{center}
\begin{tabularx}{\linewidth}{|l|X|}
\hline
\rowcolor{popblueL} \textbf{Situation} & \textbf{Résultat} \\
\hline
Équation $ax+b=0$ ($a\neq 0$) & Solution unique : $x=-\dfrac{b}{a}$, \; $\mathcal{S}=\left\{-\dfrac{b}{a}\right\}$ \\
\hline
Équation $0x+b=0$ & $b=0$ : tout réel est solution ; $b\neq 0$ : $\mathcal{S}=\varnothing$ \\
\hline
Produit nul $(ax+b)(cx+d)=0$ & $x=-\dfrac{b}{a}$ ou $x=-\dfrac{d}{c}$ \\
\hline
Ajouter/retrancher un nombre & Le sens de l'inégalité est \textbf{conservé} \\
\hline
Multiplier/diviser par $k>0$ & Le sens de l'inégalité est \textbf{conservé} \\
\hline
Multiplier/diviser par $k<0$ & Le sens de l'inégalité est \textbf{changé} \\
\hline
Signe de $ax+b$ ($a>0$) & Négatif avant $-\dfrac{b}{a}$, nul en $-\dfrac{b}{a}$, positif après \\
\hline
\end{tabularx}
\end{center}

\textbf{Méthodes express}
\begin{itemize}
\item \textbf{Équation :} développer, réduire, regrouper les $x$ d'un côté et les constantes de l'autre, puis diviser par le coefficient de $x$. \textbf{Toujours vérifier} la solution en la remplaçant dans l'énoncé.
\item \textbf{Équation produit nul :} factoriser si nécessaire, puis résoudre chaque facteur égal à zéro.
\item \textbf{Inéquation :} mêmes transformations que pour une équation, mais \textbf{changer le sens} dès qu'on multiplie ou divise par un nombre négatif. Donner $\mathcal{S}$ sous forme d'intervalle et le représenter sur la droite graduée.
\item \textbf{Problème concret :} choix de l'inconnue $\to$ mise en équation $\to$ résolution $\to$ vérification $\to$ conclusion avec une phrase et les unités.
\end{itemize}
\end{bilan}

\begin{aretenir}[Checklist avant le BEPC]
\begin{itemize}
\item[$\square$] Je sais reconnaître une équation du premier degré et la mettre sous la forme $ax+b=0$.
\item[$\square$] Je sais que la solution de $ax+b=0$ ($a\neq 0$) est $x=-\dfrac{b}{a}$.
\item[$\square$] Je sais résoudre une équation avec des parenthèses ou des fractions (développer, réduire au même dénominateur).
\item[$\square$] Je sais utiliser la règle du produit nul : $A\times B=0 \iff A=0$ ou $B=0$.
\item[$\square$] Je vérifie toujours ma solution en la remplaçant dans l'équation de départ.
\item[$\square$] Je sais qu'ajouter ou retrancher un même nombre ne change pas le sens d'une inégalité.
\item[$\square$] Je n'oublie jamais de \textbf{changer le sens} de l'inégalité quand je multiplie ou divise par un nombre strictement négatif.
\item[$\square$] Je sais donner l'ensemble des solutions d'une inéquation sous forme d'intervalle et le représenter sur une droite graduée (crochet ouvert/fermé).
\item[$\square$] Je sais dresser le tableau de signes de $ax+b$ avec la valeur qui annule : $x=-\dfrac{b}{a}$.
\item[$\square$] Je sais distinguer inégalités strictes ($<$, $>$) et larges ($\leq$, $\geq$) dans la réponse.
\item[$\square$] Je sais mettre un problème concret en équation : inconnue, équation, résolution, phrase de conclusion.
\item[$\square$] Je rédige ma démarche proprement et je justifie chaque étape.
\end{itemize}
\end{aretenir}

\findecours

\end{document}
